% concatenated from Projection-to-hyperplane.tex
\documentclass[11pt,a4paper,reqno]{amsart}
\usepackage[english]{babel}
\usepackage[T1]{fontenc}
\usepackage{verbatim}
\usepackage{palatino}
\usepackage{amsmath}
\usepackage{mathabx}
\usepackage{amssymb}
\usepackage{amsthm}
\usepackage{amsfonts}
\usepackage{graphicx}
\usepackage{esint}
\usepackage{color}
\usepackage{mathtools}
\usepackage{overpic}
\DeclarePairedDelimiter\ceil{\lceil}{\rceil}
\DeclarePairedDelimiter\floor{\lfloor}{\rfloor}
\usepackage[colorlinks = true, citecolor = black]{hyperref}
\pagestyle{headings}


\newcommand{\im}{\mathrm{Im}}
\newcommand{\R}{\mathbb{R}}
\newcommand{\W}{\mathbb{W}}
\newcommand{\He}{\mathbb{H}}
\newcommand{\N}{\mathbb{N}}
\newcommand{\Q}{\mathbb{Q}}
\newcommand{\C}{\mathbb{C}}
\newcommand{\Z}{\mathbb{Z}}
\newcommand{\tn}{\mathbb{P}}
\newcommand{\J}{{\cal J}}
\newcommand{\calT}{\mathcal{T}}
\newcommand{\calL}{\mathcal{L}}
\newcommand{\calD}{\mathcal{D}}
\newcommand{\calH}{\mathcal{H}}
\newcommand{\cG}{\mathcal{G}}
\newcommand{\calC}{\mathcal{C}}
\newcommand{\calE}{\mathcal{E}}
\newcommand{\calN}{\mathcal{N}}
\newcommand{\calB}{\mathcal{B}}
\newcommand{\calG}{\mathcal{G}}
\newcommand{\calF}{\mathcal{F}}
\newcommand{\calS}{\mathcal{S}}
\newcommand{\calQ}{\mathcal{Q}}
\newcommand{\calR}{\mathcal{R}}
\newcommand{\calU}{\mathcal{U}}
\newcommand{\calK}{\mathcal{K}}
\newcommand{\calM}{\mathcal{M}}
\newcommand{\calY}{\mathcal{Y}}
\newcommand{\calZ}{\mathcal{Z}}
\newcommand{\upor}{\overline{\operatorname{por}}}
\newcommand{\spt}{\operatorname{spt}}
\newcommand{\Hd}{\dim_{\mathrm{H}}}
\newcommand{\Pd}{\dim_{\mathrm{p}}}
\newcommand{\Bd}{\overline{\dim}_{\mathrm{B}}}
\newcommand{\calP}{\mathcal{P}}
\newcommand{\cH}{\mathcal{H}}
\newcommand{\E}{\mathbb{E}}
\newcommand{\cC}{\mathcal{C}}
\newcommand{\calA}{\mathcal{A}}
\newcommand{\1}{\mathbf{1}}
\newcommand{\Nil}{\mathrm{Nil}}
\newcommand{\spa}{\operatorname{span}}
\newcommand{\diam}{\operatorname{diam}}
\newcommand{\card}{\operatorname{card}}
\newcommand{\dist}{\operatorname{dist}}
\newcommand{\Var}{\operatorname{\mathbf{Var}}}
\newcommand{\sgn}{\operatorname{sgn}}
\newcommand{\Rea}{\operatorname{Re}}
\newcommand{\Ai}{\operatorname{Ai}}
\newcommand{\Bi}{\operatorname{Bi}}
\newcommand{\Bdim}{\overline{\dim_{\textsc{B}}}}
\newcommand{\interior}{\operatorname{int}}
\newcommand{\Om}{\Omega}
\newcommand{\len}{\mathrm{length}}
\newcommand{\Hei}{{\mathbb{H}}^{1}}
\newcommand{\modf}{{\mathrm{mod}_4}}
\newcommand{\finv}{{f^{-1}}}
\newcommand{\bmo}{\mathrm{BMO}}
\newcommand{\bmol}{{\mathrm{BMO}_{loc}}}
\newcommand{\pv}{\mathrm{p.v. }}
\newcommand{\Div}{\mathrm{div}_{\He}}
\newcommand{\G}{\mathbb{G}}
\newcommand{\V}{\mathbb{V}}
\newcommand{\vis}{\mathrm{Vis}}
\newcommand{\Tan}{\mathrm{Tan}}
\renewcommand{\v}{\mathfrak{v}}
\renewcommand{\Im}{\mathrm{Im\,}}
\newcommand{\m}{\mathfrak{m}}
\newcommand{\M}{\mathfrak{M}}
\newcommand{\w}{\mathfrak{w}}




\def\fz{\infty}
\def\az{\alpha}
\def\ac{{\mathrm ac}}
\def\dist{{\mathop\mathrm{\,dist\,}}}
\def\loc{{\mathop\mathrm{\,loc\,}}}
\def\weak{{\mathop\mathrm{ \,weak\,}}}
\def\lip{{\mathop\mathrm{\,Lip}}}
\def\llc{{\mathop\mathrm{\,LLC}}}
\def\carr{{\mathop\mathrm{\,Carrot\,}}}
\def\lz{\lambda}
\def\dz{\delta}
\def\bdz{\Delta}
\def\ez{\epsilon}
\def\kz{\kappa}
\def\bz{\beta}
\def\fai{\varphi}
\def\gz{{\gamma}}
\def\oz{{\omega}}
\def\Oz{{\Omega}}
\def \boz{{\Omega}}
\def\toz{{\wz{ U}}}
\def\ttoz{{\wz{\toz}}}
\def\vz{\varphi}
\def\tz{\theta}
\def\sz{\sigma}
\def\pa{\partial}
\def\pau{\partial{U}}
\def\pav{\partial{V}}
\def\wz{\widetilde}


\def\la{\langle}
\def\ra{\rangle}

\def\Barint_#1{\mathchoice
          {\mathop{\vrule width 6pt height 3 pt depth -2.5pt
                  \kern -8pt \intop}\nolimits_{#1}}%
          {\mathop{\vrule width 5pt height 3 pt depth -2.6pt
                  \kern -6pt \intop}\nolimits_{#1}}%
          {\mathop{\vrule width 5pt height 3 pt depth -2.6pt
                  \kern -6pt \intop}\nolimits_{#1}}%
          {\mathop{\vrule width 5pt height 3 pt depth -2.6pt
                  \kern -6pt \intop}\nolimits_{#1}}}


\numberwithin{equation}{section}

\theoremstyle{plain}
\newtheorem{thm}[equation]{Theorem}
\newtheorem*{"thm"}{"Theorem"}
\newtheorem{conjecture}[equation]{Conjecture}
\newtheorem{lemma}[equation]{Lemma}
\newtheorem{sublemma}[equation]{Sublemma}
\newtheorem{ex}[equation]{Example}
\newtheorem{cor}[equation]{Corollary}
\newtheorem{proposition}[equation]{Proposition}
\newtheorem{question}{Question}
\newtheorem{claim}[equation]{Claim}
\newtheorem{fact}[equation]{Fact}
\newtheorem{problem}{Problem}
\newtheorem*{counter}{Counter assumption}

\theoremstyle{definition}
\newtheorem{conditionC}{Condition (C')}
\newtheorem{definition}[equation]{Definition}
\newtheorem{assumption}[equation]{Assumption}
\newtheorem{notation}[equation]{Notation}
\newtheorem{observation}[equation]{Observation}

\theoremstyle{remark}
\newtheorem{remark}[equation]{Remark}
\newtheorem{terminology}[equation]{Terminology}
\newtheorem{discussion}[equation]{Discussion}
\newtheorem{assumptions}[equation]{Assumptions}
\newtheorem{warning}[equation]{Warning}

\addtolength{\hoffset}{-1.15cm}
\addtolength{\textwidth}{2.3cm}
\addtolength{\voffset}{0.45cm}
\addtolength{\textheight}{-0.9cm}

\newcommand{\nref}[1]{(\hyperref[#1]{#1})}

\DeclareMathSymbol{\intop}  {\mathop}{mathx}{"B3}
\setcounter{section}{0}


\author{Jiayin Liu}
\title[Restricted Projections to Hyperplanes in $\mathbb{R}^n$]{Restricted Projections to Hyperplanes in $\R^n$}
\address{Mathematics Area, SISSA, Trieste, Italy}
\email{jliu@sissa.it}
%\date{}
\subjclass[2010]{28A80, 28A78}
\keywords{Restricted projection, Cinematic maps, Hausdorff dimension, Line-cone incidence}
\thanks{}

\begin{document}




\begin{abstract}
We study dimensions of sets projected to an $(n-2)$-dimensional family of hyperplanes in $\mathbb{R}^n$ under curvature conditions.
Let $n\ge 3$ and $\Sigma \subset S^{n-1}$ be an $(n-2)$-dimensional $C^2$ manifold such that $\Sigma$ has non-vanishing geodesic curvature ($n=3$)/sectional curvature $>1$ ($n \ge 4)$. Let $Z \subset \mathbb{R}^{n}$ be analytic with $\dim Z  \le n-2$ and $0 < s < \dim Z$. Then
\begin{equation*}
  \dim \{x \in \Sigma : \dim \pi_{T_xS^{n-1}}(Z) < s\} \le s
\end{equation*}
where $\pi_{T_xS^{n-1}}$ is the orthogonal projection from $\mathbb{R}^n$ to the tangent space $T_xS^{n-1}$.
In particular,
for $\mathcal{H}^{n-2}$-a.e. $x \in \Sigma$, $\dim \pi_{T_xS^{n-1}}(Z) = \dim Z$. When $n=3$ and $\dim Z < 1$, the quantitative estimate improves the one obtained by Gan-Guo-Guth-Harris-Maldague-Wang.

For the case $\dim Z > n-2$, if in addition $\pi_{T_yS^{n-1}}(Z) \le n-2$ for some $y \in S^{n-1}$, we show that
$\dim \pi_{T_xS^{n-1}}(Z) = \min\{\dim Z, n-1\}$ for $\mathcal{H}^{n-2}$-a.e. $x \in \Sigma$.
%The main idea is to transfer the study to a line-cone incidence problem in $\R^n$.
\end{abstract}

\maketitle
\tableofcontents

\section{Introduction}
This paper focuses on studying projection to a restricted family of
hyperplanes for any analytic set in $\R^n$. We first fix several notations.
Denote by $\pi_W: \R^n \to W$ the orthogonal projection to the ($k$-dimensional, $k=1, \cdots, n-1$) subspace $W \subset \R^n$ and by $T_x M$ the tangent space at $x \in M$ of a differentiable manifold $M \subset\R^n $. The notation $\dim$ will always refer to Hausdorff dimension.



In \cite{FasslerOrponen14}, F\"{a}ssler and Orponen conjectured that
\begin{conjecture} \label{conjfo}
   Let $Z$ be an analytic set in $\R^3$ and $\gz \in C^2([0,1], S^2)$ be a non-degenerate curve, that is
   \begin{equation}\label{nondeg0}
    {\rm span}\{ \gamma(\theta), \dot \gamma(\theta), \ddot \gamma(\theta)\}=\mathbb{R}^3 \text{  for all }  \theta \in [0,1].
\end{equation}
    Then
    \begin{enumerate}
      \item[(i)]
 \begin{equation*}
   \dim [\pi_{{\rm span}\{\gamma(\theta)\}}(Z)]  = \min\{\dim Z, 1 \}, \quad \calH^1\mbox{-\rm a.e.} \ \theta\in [0,1];
\end{equation*}
      \item[(ii)]
      \begin{equation*}
   \dim [\pi_{T_{\gamma(\theta)}S^2}(Z)]  = \min\{\dim Z, 2 \}, \quad \calH^1\mbox{-\rm a.e.} \ \theta\in [0,1].
\end{equation*}
    \end{enumerate}
\end{conjecture}
Conjecture \ref{conjfo} (i) was confirmed by Pramanik-Yang-Zahl \cite{2022arXiv220702259P} and Conjecture \ref{conjfo} (ii) was confirmed by Gan-Guo-Guth-Harris-Maldague-Wang \cite{gan2024}.

A higher
dimensional extension of Conjecture \ref{conjfo} (i) was studied in \cite{LIU2025}, which states that
\begin{thm}\label{2025}
  For $n \ge 4$, let $Z$ be an analytic set in $\R^n$ and $\Sigma \subset S^{n-1}$ be an $(n-2)$-dimensional $C^2$ manifold with sectional curvature $>1$.
    Then
 \begin{equation*}
   \dim [\pi_{{\rm span}\{x\}}(Z)]  = \min\{\dim Z, 1 \}, \quad \calH^{n-2} \mbox{-\rm a.e.} \ x\in \Sigma.
\end{equation*}
\end{thm}

This paper focuses on higher dimensional extension of Conjecture \ref{conjfo} (ii) in the spirit of Theorem \ref{2025}.
\begin{conjecture} \label{conjl}
   For $n \ge 4$, let $Z$ be an analytic set in $\R^n$ and $\Sigma \subset S^{n-1}$ be an $(n-2)$-dimensional $C^2$ manifold with sectional curvature $>1$.
    Then
      \begin{equation*}
   \dim [\pi_{T_{x}S^{n-1}}(Z)]  = \min\{\dim Z, n-1 \}, \quad \calH^{n-2} \ a.e. \ x\in \Sigma.
\end{equation*}
\end{conjecture}

In this paper, we show that Conjecture \ref{conjl} is true if $\dim Z \le n-2$ and  we also make partial progress when $\dim Z >n-2$. Note that the non-degenerate condition \eqref{nondeg0} is equivalent to $\gamma$ having non-vanishing geodesic curvature (c.f. \cite[Page 3]{LIU2025}).

\begin{thm}\label{planemain}
For $n \ge 3$, let $Z$ be an analytic set in $\R^n$ with $\dim Z \le n-2$ and $\Sigma \subset S^{n-1}$ be an $(n-2)$-dimensional $C^2$ manifold having
\begin{itemize}
  \item non-vanishing geodesic curvature if $n=3$;
  \item sectional curvature $>1$ if $n \ge 4$.
\end{itemize}
Then, for any $s < \dim Z$, we have
  $$  \dim \{ x \in \Sigma : \dim [\pi_{T_{x}S^{n-1}}(Z)] < s  \} \le s.  $$
  In particular,
\begin{equation*}
   \dim [\pi_{T_{x}S^{n-1}}(Z)]  = \dim Z, \quad \calH^{n-2} \mbox{-\rm a.e.} \ x \in \Sigma.
\end{equation*}
\end{thm}

\begin{remark} \label{rmkmain0}
  When $n=3$, under the same assumption on $\Sigma$, Gan-Guo-Guth-Harris-Maldague-Wang \cite[Theorem 1]{gan2024} showed that for any $Z\subset \R^3$,
    $$  \dim \{ x \in \Sigma : \dim [\pi_{T_{x}S^{2}}(Z)] < s  \} \le s + 1-\dim Z.  $$
  Thus, if $\dim Z < 1$, Theorem \ref{planemain} improves their estimate.
\end{remark}

\begin{thm}\label{cormain}
For $n \ge 3$, let $Z$ be an analytic set in $\R^n$ with $\dim Z > n-2$ and $\Sigma \subset S^{n-1}$ be as in Theorem \ref{planemain}.
Moreover, if there exists $y \in S^{n-1}$ such that
   \begin{equation}\label{le1}
     \dim[\pi_{T_{y}S^{n-1}}(Z)] \le n-2,
   \end{equation}
  then,
  \begin{equation*}
   \dim [\pi_{T_{x}S^{n-1}}(Z)]  = \min\{\dim Z, n-1 \}, \quad \calH^{n-2} \mbox{-\rm a.e.} \ x \in \Sigma.
\end{equation*}
\end{thm}



\begin{remark} \label{rmkmain}
  \rm
(i) Notice that for sets $Z$ with $n-2<\dim Z \le n-1$, it is easy to see that $\dim[\pi_{W}(Z)] \ge \dim Z-1$ for all hyperplanes $W$. Since $\dim Z-1 \le n-2$, assumption \eqref{le1} is realisable.

  (ii) If one wants to fully resolve Conjecture \ref{conjl}, it suffices to further consider sets $Z$ with $n-2<\dim Z \le n-1$ such that for some $y \in S^{n-1}$,
     \begin{equation}\label{les}
     \dim[\pi_{T_{y}S^{n-1}}(Z)] \le s \mbox{ for some $s < \dim Z$}.
   \end{equation}
   Otherwise, $\dim[\pi_{T_{x}S^{n-1}}(Z)] =\dim Z$ everywhere and the conjecture holds for $Z$ trivially.
   Thus, if one could relax the condition \eqref{le1} in Theorem \ref{cormain} to \eqref{les}, one resolves Conjecture \ref{conjl}.
\end{remark}


\begin{ex} \label{eg1}
   \rm
   (i) Circles in $S^2$ that are not equators have non-vanishing geodesic curvature. Similarly, $(n-2)$-dimensional spheres that are not equatorial spheres have sectional curvature $>1$.
   In particular, let $\Sigma_c: = S^{n-1} \cap \{x_{n}=c\}$ for $c \in (0,1)\cup(-1,0)$. Indeed, it is straightforward to check that $\Sigma_c$ has geodesic curvature $\pm c/\sqrt{1-c^2} \ne 0$ (depending on the orientation of $\Sigma_c$) when $n=3$ and
   sectional curvature $1/(1-c^2) >1$ when $n \ge 4$.


   (ii) Other examples include small $C^2$-perturbations of circles and spheres in (i), since the curvature condition in Theorem \ref{planemain} is stable under small perturbations.
\end{ex}












\subsection{Background of (Restricted) Projections}
Projection properties was first established by Marstrand \cite{MR0063439} who showed that the orthogonal projection of any Borel set $Z \subset \R^2$ to almost every line has dimension $\min\{1, \dim Z\}$. For all $n \ge 2$ and $1 \le m \le n$, Mattila \cite{MR0409774} showed that the orthogonal projection of any Borel set $Z \subset \R^n$ to almost every $W \in G(n,m)$ has dimension $\min\{m, \dim Z\}$ where $G(n,m)$ is the $m(n-m)$-dimensional Grassmann manifold of all
$m$-dimensional subspaces of $\R^n$ and ``almost every'' refers to the $m(n-m)$-dimensional Hausdorff measure on $G(n,m)$.

Restricted projection focuses on studying orthogonal projection of a set in $\R^n$ to a family of $m$-subspaces which forms a submanifold $\Sigma \subset G(n,m)$. Noting that when $m=1$ or $m=n-1$, $G(n,m)$ can be naturally identified with the unit sphere $S^{n-1} \subset \R^n$. In particular, when $n=3$, $\dim G(3,1) = \dim G(3,2) = 2$, thus the only meaningful submanifold $\Sigma \subset G(3,m)$ ($m=1,2$) to consider in the restricted projection is $1$-dimensional, which is a curve in $S^2$. These cases are relatively well understood by the aforementioned Conjecture \ref{conjfo} and its resolution in \cite{2022arXiv220702259P,gan2024}, together with other previous partial progress \cite{Jarvenpaa2008,FasslerOrponen14,Jarvenpaa2014,Chen2018,MR4148151,
OrponenVenieri2020,Harris2022,Harris2023,marstrandtype}, etc.
For the case $\Sigma \subset G(n,1)$ being a $C^\infty$-curve, certain restricted projection property was established by Gan-Guo-Wang \cite{gan2022restricted}.
In general, when $2 \le m \le n-2$ and $1 \le \dim \Sigma \le \dim G(n,m)$, the study of restricted projection are widely open.


\begin{remark} \label{level} \rm
   We remark that the case $n=3$ in Theorem \ref{planemain} is also a consequence of \cite[Theorem 3.2]{Jarvenpaa2008} (see also \cite[Proposition 1.5]{FasslerOrponen14}) using the sublevel set technique pioneered by Kaufman \cite{Kaufman1968}. One could also run this technique for $n \ge 4$ in Theorem \ref{planemain} (e.g. following Page 8-9 in \cite{FasslerOrponen14}) and it however, only works for sets $Z \subset \R^n$ with $\dim Z \le 1$ instead of $\dim Z \le n-2$ in Theorem \ref{planemain}. The reason is that the proof relies on a sublevel set estimate established in \cite[Lemma 3.3]{Christ1985} for $n=3$. When $n \ge 4$, one could use the corresponding sublevel set estimate \cite[Lemma 3.4]{Christ1985} and obtain that Theorem \ref{planemain} holds for $Z \subset \R^n$ with $\dim Z \le 1$.
\end{remark}



\subsection{Ideas of the Proof of Theorem \ref{planemain}}
Let $Z \subset B(o,1/2) \subset \R^n$ be an analytic set with
$$  t:=\dim Z \in (0,n-2].   $$
By proving the theorem on each local chart of $\Sigma$, we can assume $\Sigma \subset \mathbb{R}^{n+1}$ is a $C^2$ diffeomorphism from $[0,1]^{n-2}$ to $\mathbb{R}^{n}$ and therefore throughout the paper we assume $\Sigma \in C^2([0,1]^{n-2}, \mathbb{R}^{n})$ to be a $C^2$ diffeomorphism.
Thus the differential $\nabla \Sigma(x)$ at $x$ maps $\{\partial_1, \cdots, \partial_{n-2}\}$ (standard basis of $[0,1]^{n-2}$) to a basis
$$  \{\nabla \Sigma(x)\cdot\partial_1, \cdots, \nabla \Sigma(x)\cdot\partial_{n-2}\}$$
of $T_{\Sigma(x)}\Sigma$.
Write
%$$ \{e_i=e_i(x)\}_{1 \le i \le n-2}  \mbox{ be the Gram-Schmidt orthogonalization of $\{\nabla \Sigma(x)\cdot\partial_1\}_{1 \le i \le n-2}$}$$
\begin{equation}\label{eix}
  e_i=e_i(x)=M_\Sigma^{-1}\nabla \Sigma(x)\cdot\partial_i, \quad \forall  i =1, \cdots, n-2
\end{equation}
where $M_\Sigma: = \max_{x \in \Sigma}\{\nabla \Sigma(x)\}$
and
\begin{equation}\label{nu}
  P_x:=T_{\Sigma(x)}S^{n-1}=T_{\Sigma(x)}\Sigma \oplus {\rm span} \{\nu\} = {\rm span} \{e_1, \cdots, e_{n-2}, \nu\}
\end{equation}
where $\nu$ is the unit normal of $\Sigma$ in $S^{n-1}$.
Thus, in coordinates, the projection $\pi_{P_x}$ can be written as
$$  \pi_{P_x}(z)=(e_1\cdot z, \cdots, e_{n-2}\cdot z,  \nu\cdot z), \quad z \in \R^n.$$
Indeed, $P_x$ is a distorted orthogonal projection and for any set $Z\subset \R^n$ and $\pi_{P_x}(Z)$ has the same dimension as the orthogonal projection of $Z$ onto $T_{\Sigma(x)}S^{n-1}$.

If $z \in B(o,1/2)$ is fixed, then the map
\begin{equation}\label{fw}
  [0,1]^{n-2} \ni x \mapsto f_z(x):=(e_1(x)\cdot z, \cdots, e_{n-2}(x)\cdot z,  \nu(x)\cdot z) \in [0,1]^{n-1}
\end{equation}
maps $[0,1]^{n-2}$ into $[0,1]^{n-1}$.
The graph $f_z$ of this map is an $(n-2)$-dimension manifold in $[0,1]^{2n-3}$:
$$ \mbox{graph}(f_z)(x):=(x_1, \cdots, x_{n-2},e_1(x)\cdot z, \cdots, e_{n-2}(x)\cdot z,  \nu(x)\cdot z) \quad x \in [0,1]^{n-2}.$$
For instance, if $n=3$, $\mbox{graph}(f_z)$ is a curve in $[0,1]^3$ for each $z \in Z$.
Observing that
$$ \bigcup_{x \in [0,1]^{n-2}}\pi_{P_x}(Z) = \bigcup_{x \in [0,1]^{n-2}}\bigcup_{z \in Z }\pi_{P_x}(z) = \bigcup_{z \in Z }\bigcup_{x \in [0,1]^{n-2}}(x,\pi_{P_x}(z))=\bigcup_{z \in Z}\mbox{graph}(f_z), $$
we see that the study of restricted projection can be transferred to the study of the union of graphs of $\{f_z\}_{z \in Z}$.

To be precise,
in order to show Theorem \ref{planemain}, fix $0<s<s' \le \dim Z \le n-2$ and choose any subset $X \subset [0,1]^{n-2}$ with $\calH^{s'}(X)>0$.
Then the fact that Theorem \ref{planemain} holds is equivalent to that
\begin{equation}\label{heu1}
  \dim (\pi_{P_x}(Z)) = \dim (\bigcup_{z \in Z} [(\{x\} \times \R^{n-1}) \cap \mbox{graph}(f_z)]) \ge s, \quad \mathcal{H}^{s'}\mbox{-a.e. } x \in X.
\end{equation}
For any $z \in Z$, let $X_z \subset X$ with $\calH^{s'}(X_z) \sim \calH^{s'}(X)$ and consider the set
\begin{equation}\label{heu3}
    \calE:  = \bigcup_{z \in Z} [(X_z \times [0,1]^{n-1}) \cap \mbox{graph}(f_z)]= \bigcup_{z \in Z} \bigcup_{x \in X_z}(x,f_z(x)) \subset [0,1]^{2n-3}.
\end{equation}
We will further verify \eqref{heu1} by showing that $\calE$ has dimension lower bound
\begin{equation}\label{heu2}
  \dim \calE \ge s' +s.
\end{equation}
This dimension estimate will be done by a $\dz$-discretised proof on the following estimate:
\begin{thm}\label{thmconfig}
For $0<s \le s' \le 1$ and $\epsilon>0$, there exist $\delta_{0}=\delta_{0}(\epsilon,s,s',n) \in (0,\tfrac{1}{2}]$ such that the following holds for all $\delta \in (0,\delta_{0}]$:
$$|\mathcal{E}|_{\delta} \geq \delta^{16\epsilon - s- s'}$$
where $|\cdot|_\delta$ stands for the $\delta$-covering number of a set.
\end{thm}


We briefly explain how to show Theorem \ref{thmconfig}. Given $\eta>0$ and a subdomain $D \subset [0,1]^{m}$, the vertical $\eta$-neighborhood of the graph of $f:D \to \R^l$ is
\begin{equation}\label{vertical}
  f^{\eta}(D):= \{(x,y) \in D \times \R^{l} : |(x,y)-(x,f(x))|= |y-f(x)|< \eta \} \subset \R^{m+l} .
\end{equation}
Noting that $\calE$ is a union of subsets contained in graphs of $\{f_z\}_{z \in Z}$. To get the desired lower bound,
the core of the proof of Theorem \ref{thmconfig} is to upper bound the volume of the intersection $f^\delta_z \cap f^\delta_{z'}$ of two $\delta$-vertical neighbourhoods in $[0,1]^{2n-3}$.
We have the following estimate.

\begin{lemma} \label{pairvolume2}
We have
  $$ \calL^{2n-3} (f^\delta_z([0,1]^{n-2}) \cap f^\delta_{z'}([0,1]^{n-2})) \lesssim \frac{\delta^{2n-3}}{\|f_z-f_{z'}\|_{C^2([0,1]^{n-2})}^{n-2}} $$
  where $\calL^{d}$ is the $d$-dimensional Lebsegue measure.
\end{lemma}
In fact, we will show the estimate holds for a more general family of maps and also show Theorem \ref{thmconfig} holds in a more general setting. See Lemma \ref{fg} in Section \ref{s2} and Theorem in Section \ref{s3} respectively.

The estimate in Lemma \ref{pairvolume2} shows that the intersection looks like the $\delta$-neighbourhood of the intersection of two $(n-2)$-dimensional planes intersecting at a point (e.g. span$\{x_1,x_2\}$  and span$\{x_3,x_4\}$ in $\R^5$ when $n=3$).
The proof relies on the fact that the image
\begin{equation}\label{szast}
  \Sigma^\ast := \nu([0,1]^{n-2}) \subset S^{n-1}
\end{equation}
of the normal vector field $\nu:[0,1]^{n-2} \to S^{n-1}$
is a $C^2$-manifold with non-vanishing geodesic curvature if $n=3$ or sectional curvature $>1$ if $n \ge 4$, which is a consequence of $\Sigma$ having non-vanishing geodesic curvature/sectional curvature $>1$ (see Lemma \ref{dual} in Section \ref{s2}). We will call $\Sigma^\ast$ the dual manifold of $\Sigma$.

This completes the outline of the proof of Theorem \ref{planemain}.







\subsection{Ideas of the Proof of Theorem \ref{cormain}}
Let $Z \subset B(o,1/2) \subset \R^n$ be an analytic set with
$$  t:=\dim Z \in (n-2,n-1].   $$
By the same argument as in the previous section and under the same notation, recalling \eqref{heu1} - \eqref{heu2},
the fact that Theorem \ref{cormain} holds is equivalent to that
\begin{equation}\label{heu11}
  \dim (\pi_{P_x}(Z)) = \dim (\bigcup_{z \in Z} [(\{x\} \times \R^{n-1}) \cap \mbox{graph}(f_z)])  =t, \ \mathcal{H}^{n-2}\mbox{-a.e. } x \in [0,1]^{n-2}.
\end{equation}
Again, for each $z \in Z$, let $X_z \subset [0,1]^{n-2}$ with $\calL^{n-2}(X_z) \sim 1$ and consider the set
\begin{equation}\label{heu31}
    \mathcal{Z}: = \bigcup_{z \in Z} [(X_z \times [0,1]^{n-1}) \cap \mbox{graph}(f_z)] = \bigcup_{z \in Z} \bigcup_{x \in X_z}(x,f_z(x)) \subset [0,1]^{2n-3}.
\end{equation}
and \eqref{heu11} will be verified by showing that $\calE$ has dimension
\begin{equation}\label{heu21}
  \dim \mathcal{Z} = n-2 + \dim Z = n-2 + t,
\end{equation}
which will be further guaranteed by the following lemma.
\begin{lemma}\label{volume2}
For any $\epsilon>0$, there exists $\delta_{0}=\delta_{0}(\epsilon) \in (0,\tfrac{1}{2}]$ such that the following holds for all $\delta \in (0,\delta_{0}]$:
$$|\mathcal{Z}|_{\delta} \gtrsim \delta^{3\epsilon - (n-2 +t)}.$$
\end{lemma}
We give the heuristic idea to show Lemma \ref{volume2}. We choose a $\delta$-separated subset of $Z$ (still denoted by $Z$) with cardinality $\sim \delta^{-t}$ and in the following of this section $Z$ will always refers to this discrete set .
It suffices to show
$$  \calL^{2n-3}(\bigcup_{z \in Z}f_z^\delta([0,1]^{n-2})) \gtrsim    \delta^{O(\epsilon) + 2n-3-(t+ n-2)} = \delta^{O(\epsilon) + n-1-t}.$$
However, only the volume estimate in Lemma \ref{pairvolume2} is not enough for this lower bound. We also need to control the total number of pairs $(z,z') \in Z\times Z$ such that $f_z^\delta([0,1]^{n-2}) \cap f_{z'}^\delta([0,1]^{n-2}) \ne \emptyset$. This is because when $t > n-2$, the total number of pairs in $Z\times Z$ is $\delta^{-2t} > \delta^{-2(n-2)}$ which is large.
Hence,
we will obtain Lemma \ref{pairvolume2} by establishing a line-cone incidence estimate together with the help of the assumption \eqref{le1} in Theorem \ref{cormain}. We sketch the ideas below.




Note that $f_z = f_{z'}$ at $x$, i.e. $f_z(x)=f_{z'}(x)$ is equivalent to that
$$  \pi_{P_x}(z) = \pi_{P_x}(z'),$$
which, together with $P_x^\perp = (T_{\Sigma(x)}S^{n-1})^\perp= $span$\{\Sigma(x)\}$ implies that
$$  z' = z + u_z \Sigma(x) \quad \mbox{ for some $u_z \ne 0$},$$
since all points in $\R^n$ which have same image as $z$ under $\pi_{P_x}$ are exactly the line $z + \R \Sigma(x)$.
This shows that for a fixed point $z \in Z$, all other points $z' \in Z \subset B(o,1/2)$ such that $f_z \cap f_{z'} \ne \emptyset$ are contained in
$$  \calC_z:=  \bigcup_{x \in [0,1]^{n-2}, r \in [-1,1] } z + r \Sigma(x)= : \bigcup_{x \in [0,1]^{n-2}} L_{z,x},  $$
which is a cone with apex $z$ and cross-section $\Sigma$.
If $\Sigma$ is the $(n-2)$-sphere $\Sigma_c$ in Example \ref{eg1} (i), $\calC_z$ is a circular cone in $\R^n$ with opening angle $\arccos c$.
We call the line segments $\{L_{z,x}\}_{x\in [0,1]^{n-2}}$ generatrices of $\calC_z$. Obviously, $\calC_z$ is the union of all its generatrices.
We conclude that
\begin{equation*}
  f_z \cap f_{z'} \ne \emptyset \quad \Longleftrightarrow \quad z' \in \calC_z.
\end{equation*}
Thus in the $\delta$-discretised version, we will prove the following relation: for any $z \in Z$,
\begin{equation}\label{linecone0}
  |\{z' \in Z : f_z^\delta([0,1]^{n-2}) \cap f_{z'}^\delta([0,1]^{n-2}) \ne \emptyset \}| \sim |\{z' \in Z : z' \in \calC_z^\delta \}|.
\end{equation}
where $A^\delta$ be the $\delta$-neighbourhood in $\R^n$ for any subset $A \subset \R^n$.
We have the following theorem on the incidence of a line and the cone $\calC_x$. Recall that the angle of a line and a subspace in $\R^n$ is the angle of the line and its orthogonal projection to that subspace.
\begin{thm} \label{lineconen}
Let $\Sigma$ be as in Theorem \ref{cormain}.
If a line $L \subset \R^n$ satisfies $\min_{p \in \calC_z}\{\angle(L,T_p\calC_z)\}  \ge a$ for some $a \gg \delta >0$, then
  $$\calL^{n}(\calC_z^\delta \cap L^\delta) \lesssim_a \delta^n. $$
  where $\angle(L,T_p\calC_z)$ is the Euclidean angle between $L$ and the tangent space $T_p\calC_z$.
\end{thm}
The proof of this theorem is geometric and relies on the curvature condition of $\Sigma$ in Theorem \ref{cormain}.
As a consequence of Theorem \ref{lineconen}, recalling that $Z$ is $\delta$-separated, we know that
\begin{equation}\label{less1}
  |\calC_z^\delta \cap L^\delta \cap Z|\lesssim_a 1, \quad \forall z \in Z.
\end{equation}



Finally, under the assumption \eqref{le1}, that is,
\begin{equation}\label{le11}
  s : = \dim[\pi_{T_yS^{n-1}}(Z)] \le n-2 < t,
\end{equation}
we can efficiently upper bound the quantity in \eqref{linecone0} as follows.

Let $Z_y:=\pi_{T_yS^{n-1}}(Z)$. By \eqref{le11}, we can find a maximal $\delta$-separated $(\delta,s)$-set $Z'_y \subset Z_y$ with $|Z'_y| \sim \delta^{-s}$ (c.f. Definition \ref{dst} in Section \ref{s3} for $(\delta,s)$-set). For each $w \in Z'_y$, the preimage set
$$W_w:=\pi_{T_yS^{n-1}}^{-1}(\{w\}) \cap Z$$
  is contained in the line $L_{w,y} = w+ \R y$.
Since $\{w+ \R y\}_{w \in Z'_y}$ are parallel, we deduce that
\begin{equation}
  Z \subset \bigcup_{w \in Z'_y} W_w^\delta \subset \bigcup_{w \in Z'_y} L_{w,y}^\delta.
\end{equation}
Roughly speaking, we are able to use condition \eqref{le11} to partition $Z$ into subsets contained $\{W_w^\delta\}_{w \in Z'_y}$ which further contained in an $s$-family of parallel $\delta$-tubes $\{L_{w,y}^\delta\}_{w \in Z'_y}$. This, combining with \eqref{less1} implies that
$$ |\calC_z^\delta \cap Z| \le \sum_{w \in Z'_y} |\calC_z^\delta \cap L_{w,y}^\delta \cap Z| \lesssim_a |Z'_y| \lesssim \delta^{-s} < \delta^{-(n-2)}. $$
Here in this introduction we omit the details on a technical approximation process to guarantee $\min_{p \in \calC_z}\{\angle(L,T_p\calC_z)\}  \ge a$.
Recalling \eqref{linecone0}, this estimate is sufficient to deduce Lemma \ref{volume2} and then Theorem \ref{cormain}, which completes the outline of the proof of Theorem \ref{cormain}.



%Thus for each $z \in Z$, since $Z'_y$ is $\delta$-separated, there exists a unique $w_z \in Z'_y$ such that $z \in W_{w_z}^\delta$.
%Letting $r \in [\delta,1]$ be dyadic, we deduce that
%\begin{align*}
%  \calL^n(\calC_z^\delta \cap Z^\delta) & \le \sum_{w \in Z'_y}  \calL^n(\calC_z^\delta \cap W_w^\delta) \le \sum_{r =\delta}^1\sum_{w \in Z'_y, r\le|w-z| \le 2r}  \calL^n(\calC_z^\delta \cap W_w^\delta)+ \calL^n(\calC_z^\delta \cap W_{w_z}^\delta) \\
%   & \lesssim \frac{r^s}{\delta^s} \calL^n(\calC_z^\delta \cap L_z^\delta)+\calL^n(L_{w_z}^\delta \cap Z^\delta) \\
%   & \lesssim \frac{r^s}{\delta^s}\frac{\delta^n}{r} +\delta^2 \\
%   & \lesssim \delta^{-1}\delta^n = \delta^2
%\end{align*}
%where $Z^\delta$ is the $\delta$-neighbourhood of $Z$ and
%\begin{itemize}
%        \item for the first term in the first $\lesssim$, we note that by $Z_\tau$ is a $(\delta,s)$-set, $|\{z \in Z_\tau, r\le|z-x| \le 2r\}| \le \frac{r^s}{\delta^s}$ and $W_z^\delta \subset L_z^\delta$;
%        \item for the second term in the first $\lesssim$, we note that $L_{z_x} \subset \calC_x$ and $W_{z_x} \subset L_{z_x} \cap Z$;
%        \item for the first term in the second $\lesssim$, we use Theorem \ref{lineconen} and we will prove that $d(x,L_z) \sim |x-z| \sim r$;
%        \item for the second term in the second $\lesssim$, since $L_{z_x}$ is a line and $Z$ is bounded, $|L_{z_x}^\delta \cap Z^\delta|$ is contained a $\delta \times \delta \times 1$-tube;
%        \item for the last $\lesssim$, we use the assumption \eqref{le1}, i.e. $s \le 1$.
%      \end{itemize}
%As a consequence,




\bigskip
\noindent
{\bf Organization.}
Section \ref{s2} is mainly designed to show Lemma \ref{pairvolume2}.
In Section \ref{s3}, we show Theorem \ref{planemain}.
Section \ref{s4} is devoted to showing Theorem \ref{lineconen}.
Finally, we complete the paper by proving Theorem \ref{cormain} in Section \ref{s5}.

\bigskip
\noindent
{\bf Notation.}
The notation $A \lesssim B$ means that there exists a constant $C \geq 1$ such that $A \leq CB$. The two-sided inequality $A \lesssim B \lesssim A$ is abbreviated to $A \sim B$. The constant $C$ can depend on the dimension $n$, and if we fix a cinematic family $F$ of cinematic constant $K$ and doubling constant $D$, $C$ may depend on $K$ and $D$.
If we want to stress the dependence of the constant $C$ with some  parameter "$\theta$", we indicate this by writing $A \lesssim_{\theta} B$. The constant $C$ may change from line to line in the proof. For simplicity, we write $\log=\log_2$ throughout the paper.

By the notation $|A|$ we denote the cardinality of $A$ if $A$ is a set and we denote the Euclidean norm if $A$ is a point. By $\mathcal{H}^{s}(A)$ (resp. $\mathcal{L}^{k}(A)$) the $s$-dimensional Hausdorff measure (resp. $k$-dimensional Lebesgue measure) of $A$. For $\delta>0$, $A^\delta$ stands for the $\delta$-neighbourhood and $|A|_\delta$ stands for the $\delta$-covering number of $A$. For a set $A \subset \R^n$ and a point $p \in \R^n$, $A \pm \{p\}$ is the set $\{x \pm p: x \in A \}$.



%when $A$ is a finite set and the $n$-dim Lebesque measure of $A$ when $A \subset \mathbb{R}^n$ is infinite.
%For $\delta \in (0,1]$, the notation $A \lessapprox_{\delta} B$ means
%that there exists a constant $C \geq 1$ such that
%\begin{displaymath}
%A \leq C \cdot\left(1+
%\log\left(\tfrac{1}{\delta}\right)^{C}\right) B.
%\end{displaymath}
%We write $A \thickapprox_{\delta} B$ if simultaneously $A
%\lessapprox_{\delta} B$ and $B \lessapprox_{\delta} A$ hold true. If the constant $C$ is allowed to depend on a parameter "$\theta$", we indicate this by writing $A \lessapprox_{\delta,\theta} B$.

%For two vectors $u,v \in \mathbb{R}^n$,
%$u \cdot v$ and $\langle u ,v \rangle$ both denote the Euclidean inner product. Given $f \in C^2(\Omega)$ for some $k$-dim $C^2$ manifold $\Omega \subset \mathbb{R}^n$,
%$\nabla f$ and $\nabla^2 f$ denote the gradient and Hessian of $f$ respectively. Moreover, as usual $\nabla_u f:= \langle \nabla f ,u \rangle$ and $\nabla_v\nabla_u f:=  \langle \nabla^2 f \cdot u ,v \rangle$.
%If $\Omega \subset \mathbb{R}$, we write $f'=\nabla f $ and $f''=\nabla^2 f$. As usual, $T_x\Omega=\mathbb{R}^k$ is the tangent space of $x \in \Omega$ and for $\xi \in T_x\Omega$ with $|\xi|=1$, we always write $\xi \in S^{k-1}$ for simplicity.


\bigskip
\noindent
{\bf Acknowledgement.}
The author would like to thank Tuomas Orponen for his interest in this paper, in particular, for pointing out the sublevel set technique in Remark \ref{level} and providing the references  \cite{Jarvenpaa2008,FasslerOrponen14}.
%The author is also grateful to the anonymous referee for several significant corrections.




\section{Cinematic Maps and Their Intersections} \label{s2}
We define a class of cinematic maps and show that the intersection of any two cinematic maps
satisfies the upper bound in Lemma \ref{pairvolume2}. Then, we check that the maps $\{f_z\}_{z \in Z}$ arising from the restricted projection are cinematic maps.

\subsection{Cinematic Maps}
\begin{definition}[Cinematic Map]\label{cine}
  Let $F \subset C^2(U,R^l)$ where $U \subset \mathbb{R}^k$ is a domain, $k,l \in \mathbb{N}$ with $l \ge k$. We say $F$ is a cinematic family of
maps, with cinematic constant $K$ and doubling constant $D$ if the following conditions
hold:
\begin{enumerate}
  \item $F$ is contained in a ball of diameter $K$ under the usual metric on $C^2(U,R^l)$;
  \item $F$ is a doubling metric space, with doubling constant at most $D$;
  \item For all $f, g \in F$, we have
  $$  \inf_{(x,\xi) \in U \times S^{k-1}} \{ |f(x) - g(x)| + |\nabla_\xi f(x) -\nabla_\xi g(x)| \} \ge K^{-1} \|f-g\|_{C^2(U)}  $$
  where $\nabla_\xi h= \nabla h \cdot \xi$ for $h=f,g$, noting that $\nabla h$ is a linear map from $\R^{k}$ to $\R^l$ and $\xi \in S^{k-1} \subset \R^{k}$ is considered as a direction of $\R^{k}$.
%  \item The family $\{\nabla_\xi\nabla_\zeta f \}_{f \in F, \xi,\zeta \in S^{k-1}} \subset C^0(U,\R^l)$ is equicontinuous.
\end{enumerate}
\end{definition}
\begin{remark} \rm
We make a remark about (3) in Definition \ref{cine}. Let us  assume $|f(x) - g(x)|=0$ at some $x$. Then (3) implies that
$$\inf_{\xi \in  S^{k-1}}|\nabla_\xi f(x) -\nabla_\xi g(x)| \sim_K \|f-g\|_{C^1(U)}.$$
This condition guarantees that whenever $f(x)=g(x)$, the smallest singular value of $\nabla(f-g)(x)$ is not zero, i.e. $\nabla(f-g)(x)$ as an $(l \times k)$-matrix has full rank.
\end{remark}




\begin{lemma} \label{local}
Assume $\eta>0$, $n \ge 3$.
Let
$F \subset C^{2}([0,1]^{n-2},[0,1]^{n-1})$ be a cinematic family with cinematic constant $K$ and $f,g \in F$.
  %Then there exists a constant $D=D(\eta,n)$ such that the following holds:
  If $U \subset [0,1]^{n-2}$ is a (closed) dyadic cube such that its diameter is smaller than $\frac\eta{2K}$, and $k$ is one of the functions $|f-g|, |\nabla_\xi (f-g)|$ for some $\xi \in S^{n-3}$, then $k$ satisfies either one of the following two inequalities:
   \begin{itemize}
     \item[(S)] $ k(x) < \eta\|f-g\|_{C^2([0,1]^{n-2}) } \mbox{ for all } x \in U$.
     \item[(L)] $k(x) \ge \eta \|f-g\|_{C^2([0,1]^{n-2})}/2 \mbox{ for all } x \in U$.
   \end{itemize}
\end{lemma}

\begin{proof}
 Write $d=\|f-g\|_{C^2( [0,1]^{n-2}) }$.
 % $h=f-g$ and let $B$ be the closed unit ball in $C^2( U)$. Then $h \in B$. In particular, Lip$h = \|\nabla h\|_{L^\infty( U) }\le 1$.
% Since ${\rm diam \, } U < \alpha(\eta/2)$, for any $x,x' \in U$, $|x-x'| < \alpha(\eta/2)$.
% By Definition \ref{cine}, we know
% that $|\nabla_\xi \nabla_\xi (f-g)(x)-\nabla_\xi \nabla_\xi (f-g)(x')| <\eta/2$.
 Since $F$ is contained in a ball of diameter $K$, we know
 $t \le K$. In particular, we know Lip$(f-g)(x) \le K$ and Lip$|\nabla(f-g)|(x)\le K$ for all $x \in U$. We have
 $$  |k(x)-k(x')| \le \max_{y\in U}\mbox{{\rm Lip} } k(y)\cdot|x-x'| \le K\diam U < \frac{\eta}2, \quad \forall x,x' \in U,  $$
 where $k$ is one of the  functions $|f-g|, |\nabla_\xi (f-g)|$.

 Next,  suppose that the
alternative (L) fails for $k = f - g$. So there is $x_0 \in U$ such that $|k(x_0)| < \eta d/2$. And for any $x \in U$, since $|x-x_0| < \eta/(2K)$, by defining $l= k/d$, we deduce that Lip$_{[0,1]^{n-2}}l \le 1$ and
$$ \frac{|k(x)|}{d} \le \frac{|k(x_0)|}{d}+ \left| \frac{|k(x)|}{d} - \frac{|k(x_0)|}{d}\right| < \frac{\eta}{2} + |l(x)-l(x_0)| \le \frac{\eta}{2} +\mbox{ [Lip$_{[0,1]^{n-2}}l$}] |x-x_0|  \le \eta.$$
 Thus the alternative (S) holds for $k$. The proof of the case $k = |\nabla_\xi (f-g)|$ is the
same.
\end{proof}


The above lemma gives the following corollary.


\begin{cor} \label{uniq}
  Let
$F \subset C^{2}([0,1]^{n-2},[0,1]^{n-1})$ be a cinematic family of maps with cinematic constant $K$, $f,g \in F$ and $h:=f-g$. Then for each dyadic cube $U \subset [0,1]^{n-2}$ with ${\rm diam \, } U < \frac{1}{4K^2}$,

{\rm (A):} at least one of the following two conditions hold:
\begin{enumerate}
  \item[(i)]
$|h(x)| \ge \frac{1}{2K} \|h\|_{C^2([0,1]^{n-2}) }$ for all $x \in  U$.
  \item[(ii)]
$|\nabla_\xi h(x)| \ge \frac{1}{2K} \|h\|_{C^2( [0,1]^{n-2}) }$ for all $x \in  U$ and $\xi \in S^{n-3}$.
\end{enumerate}

{\rm (B):} Moreover, if (A)(ii) holds, then
\begin{equation}\label{hxy}
  |h(x)|+| h(y)| \ge \frac{\|h\|_{C^2( [0,1]^{n-2}) }}{4K}|x-y|, \quad \forall x,y \in U.
\end{equation}
\end{cor}



\begin{proof}
We first show (A):
   We apply Lemma \ref{local} to $h$ with $\eta = 1/(2K)$. As a result, we know for each dyadic cube $U \subset [0,1]^{n-2}$ with diameter ${\rm diam \, } U < \frac{1}{4K^2} $,
 either (S) or (L) holds for $k$ and all $x \in U$, where $k$ is one of the
functions $h$, or $\nabla_\xi h$. On the other hand,
since $F$ is a cinematic family of functions, if the alternative (L) does not hold for the functions $h$, then by Definition \ref{cine} (3), for all $\xi \in S^{n-2}$, (L) holds for $\nabla_\xi h=\nabla_\xi (f-g)$.

We show (B). Fix $x,y \in U$. Note that if $x=y$, \eqref{hxy} holds trivially. We thus assume
$x \ne y$ and write
 $$ y = x+ \xi_y |y-x|, \mbox{ where } \xi_y = \frac{y-x}{|y-x|} \in S^{n-3}$$
 By Definition \ref{cine} (1), we know that
$$d:=\|h\|_{C^2([0,1]^{n-2})} \le 2K < \infty$$
 and hence for any $s \in [0,|x-y|]$, we have
$$  |\nabla_{\xi_y}h(x+s\xi_y) - \nabla_{\xi_y}h(x)| \le {\rm Lip}(\nabla_{\xi_y} h)\cdot |x+s\xi_y-x| \le  ds.   $$
Let
$$ v(s):=\nabla_{\xi_y}h(x+s\xi_y) - \nabla_{\xi_y}h(x), s \in [0,|x-y|]. $$
We deduce that
\begin{equation}\label{vs}
 |v(s)| \le ds \le d|x-y|, \quad s \in [0,|x-y|].
\end{equation}
We compute
\begin{align*}
  |h(x)|+| h(y)|  & \ge |h(y)-h(x)| = |\int_0^{|x-y|}\nabla_{\xi_y}h(x+s\xi_y) \, ds | \\
   & =|\int_0^{|x-y|}\nabla_{\xi_y}h(x) \, ds + \int_0^{|x-y|}v(s) \, ds |  \\
   &\ge |\nabla_{\xi_y}h(x)||x-y| - |\int_0^{|x-y|}v(s) \, ds | .
\end{align*}
Since $\nabla_{\xi_y}h$ satisfies (A)(ii), we have $|\nabla_{\xi_y}h(x)| \ge \frac{d}{2K}$. Combining \eqref{vs}, we deduce that
$$ |h(x)|+| h(y)| \ge \frac{d}{2K}|x-y| - d|x-y|^2 = d|x-y|(\frac{1}{2K} - |x-y|).  $$
Since $\diam U < \frac{1}{4K^2}$, we know $\frac{1}{2K} - |x-y| > \frac{1}{4K}$, which gives \eqref{hxy} as desired.

  The proof is complete.
\end{proof}

Recall from \eqref{vertical} that for arbitrary $\eta>0$ and
open set $D \subset [0,1]^{n-2}$, the vertical $\eta$-neighborhood of the graph of $f:D \to \R^l$ is
   $$f^{\eta}(D):= \{(x,y) \in D \times \R^{l} : |(x,y)-(x,f(x))|= |y-f(x)|< \eta \}.$$


\begin{lemma}\label{nbhd}
  If $f \in C^1(D , \R^l)$, then
   $$\calL^{n-2+l}(f^{\eta}(D)) \lesssim \eta^{l}\cdot \|f\|_{C^1(D)}\cdot  \calL^{n-2}(D).$$
\end{lemma}
\begin{proof}
  Note that
  $$ f^{\eta}(D) = \bigcup_{|\eta'|<\eta} (D,f(D)+\eta'). $$
 Write $f=(f_1, \cdots, f_l)$ where $f_i \in C^1(D)$, $i=1,\cdots,l$.
 We furthermore have
 $$ f^{\eta}(D) = \bigcup_{|\eta'|<\eta} (D,f_1(D)+\eta', \cdots, f_l(D)+\eta'). $$
%  Note that if $(x,y)=(x,y_1,cdots,y_l)\in D \times \R^{l}$ is contained in $f^{\eta}(D)$, then necessarily we have
%  $$  |y_i - f_i(x)| < \eta.  $$
  Moreover, since the set $(D,f(D))$ has $(n-2)$-Hausdorff measure $\calH^{n-2}((D,f(D))) \le \|f\|_{C^1(D)}\calL^{n-2}(D)$, we deduce that
  $$ \calL^{n-2+l}(f^{\eta}(D)) \lesssim \eta^{l}\cdot \|f\|_{C^1(D)}\cdot  \calL^{n-2}(D), $$
  which completes the proof.
\end{proof}


\begin{lemma} \label{fg}
Let $\delta>0$,
$F \subset C^{2}([0,1]^{n-2},[0,1]^{n-1})$ be a cinematic family with cinematic constant $K$. Then, for any $f,g \in F$ with $d:=\|f-g\|_{C^2([0,1]^{n-2})} \ge \delta$, we have
\begin{equation}\label{meas}
  \calL^{2n-3}[f^{\delta}([0,1]^{n-2})\cap g^{\delta}([0,1]^{n-2})] \lesssim \frac{\delta^{2n-3}}{\|f-g\|_{C^2([0,1]^{n-2})}^{n-2}}.
\end{equation}
\end{lemma}


\begin{proof}
By
decomposing $[0,1]^{n-2}$ into dyadic cubes $\{U\}$ with with ${\rm diam \, } U < \frac{1}{4K^2}$ to prove \eqref{meas}, it suffices to show
$$ \calL^{2n-3}(f^\delta(U) \cap g^\delta(U)) \lesssim \frac{\delta^{2n-3}}{d^{n-2}}   $$
for each $U \subset[0,1]^{n-2}$.
In the following, we fix $U$ and write $f^\delta \cap g^\delta = f^\delta(U) \cap g^\delta(U)$ for simplicity.

  If $f^\delta \cap g^\delta = \emptyset$, then there is nothing to prove. Also, since $d \in [\delta,1]$, if $\delta \le d \le 4K\delta$,
   then using Lemma \ref{nbhd} and noting that $\|f\|_{C^1(U)} \le K$, we know that
   $$\calL^{2n-3}(f^\delta \cap g^\delta) \le \calL^{2n-3}(f^\delta) \sim \delta^{n-1} \sim \frac{\delta^{n-1}d^{n-2}}{d^{n-2}} \sim \frac{\delta^{2n-3}}{d^{n-2}} .$$
Thus in the following, we assume
\begin{equation}\label{ass1}
  d > 4K\delta.
\end{equation}



  Assume $f^\delta \cap g^\delta \ne \emptyset$.
  Let $P_{f,g}^U$ be the orthogonal projection of
  $f^\delta \cap g^\delta$ from $U \times [0,1]^{n-1}$ to $U$.
  We have
  $$ f^\delta \cap g^\delta \subset f^\delta(P_{f,g}^U). $$
  Thus
 by Lemma \ref{nbhd},
 we need to show
  \begin{equation}\label{Pfg}
   \mathcal{L}^{n-2} (P_{f,g}^U) \lesssim \frac{\delta^{n-2}}{d^{n-2}}.
  \end{equation}


  Write $h=f-g$. By $f^\delta \cap g^\delta \ne \emptyset$, we know there exists $x \in U$ such that $|h(x)|\le \delta < \frac{d}{4K}$ which implies $|h|$ does not enjoy Corollary \ref{uniq} (i).
  Hence Corollary \ref{uniq} (A)(ii) holds and
  we can apply Corollary \eqref{uniq} (B) to deduce that
  $$ |h(y)| \ge \frac{d}{4K} |x-y| - |h(x)| \ge \frac{d}{4K} |x-y| -\delta.  $$
  Thus if $|x-y| > 16K^2\delta/t$,
  \begin{equation}\label{hy}
    |h(y)| = |f(y)-g(y)| \ge 4K\delta.
  \end{equation}
  We claim that if $|x-y| > 16K^2\delta/d$, then $(y,g(y)) \notin f^\delta$, that is, for all $z \in U$ we have
  \begin{equation}\label{cl1}
    |(y,g(y)) - (z,f(z))| > \delta.
  \end{equation}
  Assume the claim holds. Then we can infer that
  $$   y \notin P_{f,g}^U, \mbox{ if }|x-y| > 16K^2\delta/d.  $$
  As a consequence
  \begin{equation}\label{ball}
    P_{f,g}^U \subset B(x,16K^2\delta/d ) \subset U \subset [0,1]^{n-2},
  \end{equation}
  which gives \eqref{Pfg} as desired.

  We are left to show \eqref{cl1}.
  Noting that if $|z-y|> \delta$, then \eqref{cl1} can not hold.
  Thus it only suffices to consider $z \in U$ with $|z-y| \le \delta$.

  Since $\|f\|_{C^2([0,1]^{n-2})} \le K$, we know
  $$  |f(z)-f(z')| \le K|z-z'|, \quad z,z' \in U. $$
  Thus for all $z \in U$ with $|z-y| \le \delta$, recalling \eqref{hy}, we have
  $$   |g(y)-f(z)| \ge |g(y)-f(y)| - |f(z)-f(y)| \ge 3K\delta,  $$
   which implies that
  $$  |(y, g(y))-(z, f(z)) | \ge |g(y)-f(z)|\ge 3K\delta, \quad z \in B(y,\delta) \cap U.   $$
 Thus \eqref{cl1} holds and
    the proof of Lemma \ref{fg} is complete.
\end{proof}



\subsection{$\{f_z\}_{z \in Z}$ is a Cinematic Family}

Recall the unit normal $\nu$ defined in \eqref{nu} and the dual surface $\Sigma^\ast= \nu([0,1]^{n-2}) \subset S^{n-1}$ in \eqref{szast}.


\begin{ex} \label{eg2}
   \rm
  Recall from Example \ref{eg1} (i) that
  $$ \Sigma_c:= S^{n-1} \cap \{(x_1, \cdots, x_n) : x_n =c\}, \quad c\in (-1,0)\cup(0,1). $$
  Then one could easily check that the dual $\Sigma^\ast_c$ of $\Sigma_c$ is
  $$ \Sigma_c^\ast= \Sigma_{\sqrt{1-c^2}}. $$
  We also note that $(\Sigma_c^\ast)^\ast = \Sigma_c$.
\end{ex}

We study the geometric property of the dual $\Sigma^\ast$.

\begin{lemma} \label{dual}
   Let $\Sigma \subset S^{n-1}$ be an $(n-2)$-dimensional $C^2$ manifold having
\begin{itemize}
  \item[(i)] non-vanishing geodesic curvature if $n=3$;
  \item[(ii)] sectional curvature $>1$ if $n \ge 4$.
\end{itemize}
  Then, $\Sigma^\ast$ also satisfies (i) and (ii) respectively.

  Moreover,
  \begin{enumerate}
  \item[(a)] if $\kappa_1$, $\cdots$, $\kappa_{n-2}$ are the geodesic curvature ($n=3$) or principal curvatures ($n\ge4$) at $x \in \Sigma$, then $\kappa_1^{-1}$, $\cdots$, $\kappa_{n-2}^{-1}$ are the geodesic curvature/principal curvatures at $\nu(x) \in \Sigma^\ast$;
  \item[(b)] for each $x \in \Sigma$, we have
  $$  T_{\nu(x)}\Sigma^\ast = T_{x}\Sigma  \mbox{ and } T_{\nu(x)}S^{n-1} = T_{\nu(x)}\Sigma^\ast \oplus {\rm span}\{x\}.  $$
  \end{enumerate}
\end{lemma}

\begin{proof}
   If $n=3$, the lemma is already known in \cite[Page 9]{gan2024}. Thus we only prove the case $n\ge 4$ and for $(a),(b)$, the case $n=3$ is an easier modification of the following proof. Since $\Sigma \subset S^{n-1}$, Gauss' equation reads as
   $$ K_{\xi_i,\xi_j}(x) = \kappa_{i}(x) \kappa_{j}(x) + 1 ,\quad x \in \Sigma  $$
  where $\{\xi_i\}_{1 \le i\le n-2} \subset T_x\Sigma$ are the $n-2$ principal directions,
  $K_{e_i,e_j}(x)$ is the sectional curvature of $\Sigma$ at $x$ spanned by $\xi_i$ and $\xi_j$ and $\kappa_{i}(x)$ is the principal curvature in the direction $\xi_i$. Thus the sectional curvature $>1$ is equivalent to that the $n-2$ principal curvatures at each $x \in \Sigma$, $\kappa_1(x), \cdots, \kappa_{n-2}(x)$ are non-vanishing and have same sign. Fixing an orientation of the normal vector $\nu$, we assume $\kappa_i(x)>0$ for all $x \in \Sigma$ and $i=1, \cdots, n-2$.

  It suffices to check that principal curvatures $\kappa_i^\ast$ of the dual manifold $\Sigma^\ast$ are positive at all $y \in \Sigma^\ast$.

  For each $y \in \Sigma^\ast$, $y=\nu(x)$ for some $x \in \Sigma$.
  The tangent space $T_y\Sigma^\ast$ of $\Sigma^\ast$ at $y$ is the space $\nabla \nu(T_x\Sigma)$. Since $\nabla \nu$ is the Gauss map, $\nabla \nu$ is also the second fundamental form of $\Sigma$ and thus has real eigenvalues, which are precisely the principal curvatures $\kappa_i$, $i=1,\cdots ,n-2$. Recalling that the principal directions
  $\{\xi_i=\xi_i(x)\}_{i=1, \cdots,n-2}$ are the eigenvectors of $\nabla \nu(x)$, we have
  \begin{equation}\label{eiast}
    \xi_i^\ast:= \nabla \nu(\xi_i) = \kappa_i \xi_i, \quad i=1, \cdots,n-2.
  \end{equation}
  Thus $\{\xi_i^\ast\}_{i=1, \cdots,n-2}$ also forms a basis of $T_y\Sigma^\ast$.
  Noting that
  $$ \R^n= {\rm span} \{x, \xi_1, \cdots, \xi_{n-2}, \nu(x)=y\} ={\rm span} \{x, \nu(x)=y\} \oplus T_y\Sigma^\ast,  $$
  we deduce that
  the normal vector $\nu^\ast(y)$ which satisfies
   $$\nu^\ast(y)\subset T_{y} S^{n-1} \mbox{ and } \nu^\ast(y) \perp T_{y}\Sigma^\ast$$
   is $x$, i.e. $\nu^\ast(y) = \nu^\ast(\nu(x)) = x$ for all $y = \nu(x)$, $x \in \Sigma$. Thus (b) holds.
   In particular, $\nu^\ast \circ \nu =$Id and $\nabla(\nu^\ast \circ \nu) = \nabla $Id $=$ Id where Id is the identity map.

   We then compute the principal curvatures of $\Sigma^\ast$ at $y$, which are the eigenvalues of $\nabla\nu^\ast(y)$.
   Recalling \eqref{eiast}, we have
   $$ \nabla\nu^\ast(y)(\xi_i^\ast) = \nabla\nu^\ast(\nu(x))(\xi_i) =\xi_i = \frac{1}{\kappa_i(x)}\xi_i^\ast, \quad i=1, \cdots,n-2.$$
   As a result,
   the principal curvatures of $\Sigma^\ast$ at $y$ are $\{\kappa_i^{-1}(x)\}_{i=1, \cdots,n-2}$. Thus (a) holds and the sectional curvature spanned by $\xi_i^\ast$ and $\xi_j^\ast$ is
   $$K_{\xi_i^\ast,\xi_j^\ast}(y) = \frac1{\kappa_{i}(x) \kappa_{j}(x) }+ 1 >1 ,\quad y=\nu(x), x \in \Sigma. $$
   The proof is complete.
\end{proof}

In the following of the paper, we let
$$ \underline\kappa:= \min_{x \in \Sigma,i =1,\cdots, n-2}\{\kappa_i(x)\} \mbox{ and } \bar\kappa:= \max_{x \in \Sigma,i =1,\cdots, n-2}\{\kappa_i(x)\} $$
where $\kappa_i(x)$ are the principle curvatures of $\Sigma$ at $x$.
Under the assumption of Theorem \ref{planemain}, we know
$$\bar\kappa \ge \underline\kappa >0.$$
By Lemma \ref{dual}, we know that principal curvatures $\kappa_i^\ast$ of the dual manifold $\Sigma^\ast$ satisfy
\begin{equation}\label{bdds}
  \frac1{\bar\kappa}= \min_{x \in \Sigma^\ast,i =1,\cdots, n-2}\{\kappa_i^\ast(x)\} \mbox{ and } \frac1{\underline\kappa}= \max_{x \in \Sigma,i =1,\cdots, n-2}\{\kappa_i^\ast(x)\}.
\end{equation}

\begin{remark} \label{remII}
    Recalling that the second fundamental form of the normals $\nu$ of $\Sigma$ and $\nu^\ast$ of $\Sigma^\ast$ are
    $$ T\Sigma^\ast \ni \zeta \mapsto (\nabla_\zeta \nu \cdot \zeta)(x), \ x \in \Sigma \mbox{ and } T\Sigma^\ast \ni \zeta \mapsto (\nabla_\zeta \nu^\ast \cdot \zeta)(x), \ x\in \Sigma^\ast$$
    and the eigenvalues of the second fundamental form are precisely principal curvatures,
    we deduce from \eqref{bdds} that
    \begin{equation}\label{II}
    (\nabla_\zeta \nu \cdot \zeta)(x) \le \bar \kappa|\zeta|^2, \quad \forall x \in \Sigma, \ \forall \zeta \in T_x\Sigma
\end{equation}
and
\begin{equation}\label{IIast}
    (\nabla_\zeta \nu^\ast \cdot \zeta)(x) \ge \frac{1}{\bar \kappa}|\zeta|^2, \quad \forall x \in \Sigma^\ast, \ \forall \zeta \in T_x\Sigma^\ast.
\end{equation}
\end{remark}


\begin{remark}\label{basis}
Recall from \eqref{eix} that for each $x \in [0,1]^{n-2}$, $\{e_i(x)= \nabla\Sigma(x)\cdot\partial_i\}_{1 \le i \le n-2}$ is a basis of $T_x\Sigma$.
   Since $\Sigma:[0,1]^{n-2} \to \Sigma([0,1]^{n-2})$ is a diffeomorphism, for each $x \in [0,1]^{n-2}$, setting $w=\Sigma(x)$, we define
$$  \{\tilde e_i(w):= e_i(\Sigma^{-1}(w))= \nabla \Sigma(x) \cdot \partial_i\}_{1\le i \le n-2}. $$
Recalling Lemma \ref{dual} (b), we know that
\begin{equation*}
  {\rm span}  \{w=\nu^\ast, \tilde e_1, \cdots, \tilde e_{n-2}, \nu\}  = \mathbb{R}^{n}
\end{equation*}
where $\nu$ (resp. $\nu^\ast$) is the unit normal of $\Sigma \subset S^{n-1}$ (resp. $\Sigma^\ast \subset S^{n-1}$). In other words, $\{\nu^\ast, \tilde e_1, \cdots, \tilde e_{n-2}, \nu\}$ forms a basis of $\R^n$. Thus
we know for any $p \in \R^n$, $p$ can be uniquely written as
$$  p =  a_0 \nu^\ast + a_1 \tilde e_1 + \cdots + a_{n-2} \tilde e_{n-2} + a_{n-1} \nu
$$
where $a_i\in \R$ for $i=0, \cdots, n-1$. However, since the basis $\{\nu^\ast, \tilde e_1, \cdots, \tilde e_{n-2}, \nu\}$ may not be orthonormal, we have the following lower bound
$$  \sum_{i=0}^{n-1} a_i^2 \ge C_\Sigma |p|^2  $$
for some $0<C_\Sigma \le 1$ only depending on $\Sigma$.
\end{remark}


\begin{lemma} \label{cineprop}
  Let $Z$ be a subset in $B(o,1/2)\subset \mathbb{R}^{n}$ and $\Sigma : [0,1]^{n-2} \to \mathbb{S}^{n-1}$ be as in Theorem \ref{planemain}. Then  $F:=\{f_{z} \in C^2( [0,1]^{n-2},[0,1]^{n-1})\}_{z \in Z}$ defined in \eqref{fw} is a cinematic family with the same doubling constant as $\mathbb{R}^{n}$. In particular, $F$ is biLipschitz homeomorphic to $Z$.
\end{lemma}

\begin{proof}
First we show (1) in Definition \ref{cine}.
Since
\begin{equation}\label{cine1}
  \sup_{f_{z} \in F }\|f_{z}\|_{C^2([0,1]^{n-2})} \le \|\Sigma\|_{C^2([0,1]^{n-2})} \max_{z \in Z}|z| \lesssim 1
\end{equation}
 where the implicit constant only depends on $\Sigma$, (1) holds.

Fix $y,z \in Z$. We show (3) in Definition \ref{cine}. Similar as \eqref{cine1}, we have
\begin{equation}\label{dou1}
  \|f_y-f_{z}\|_{C^1([0,1]^{n-2})} \lesssim |y-z|.
\end{equation}
Write
$h:= f_y-f_{z}$. We prove (3) by showing
\begin{equation}\label{cine2}
  \inf_{x \in [0,1]^{n-2}, \xi \in S^{n-3}} \{ |h(x)| + |\nabla_\xi h(x)|  \} \gtrsim  |y-z|.
\end{equation}
We recall from Remark \ref{basis} that $w=\Sigma(x)$ and
$$  \{\tilde e_i(w):= e_i(\Sigma^{-1}(w))= \nabla \Sigma(x) \cdot \partial_i\}_{1\le i \le n-2} \mbox{ is a basis of $T_{w}\Sigma$}. $$
Define
$\tilde h  \in C^2(\Sigma)$ by
\begin{equation}\label{th}
  \tilde h(w):= ((y-z) \cdot \tilde e_1(w), \cdots,(y-z) \cdot \tilde e_{n-2}(w), \nu(w))  \quad w \in \Sigma.
\end{equation}
Thus $h=\tilde h \circ \Sigma$. For each $\xi \in S^{n-3}$, let $\xi_\ast:= \nabla \Sigma(x) \cdot \xi \in T_{\Sigma(x)} \Sigma$. We know for all $x \in [0,1]^{n-2}$ and $\xi \in S^{n-3}$,
\begin{equation}\label{rel3}
 \nabla_\xi h (x)= \langle\nabla h (x), \xi \rangle = \langle\nabla \tilde h (\Sigma(x)), \xi_\ast \rangle = \nabla_{\xi_\ast} \tilde h (w).
\end{equation}
and in particular,
\begin{equation} \label{chain}
  \nabla_{\partial_i} h (x) = \nabla_{\tilde e_i} \tilde h (w), \quad i =1, \cdots, n-2.
\end{equation}
Noting that $|\nabla \Sigma|$ is nowhere vanishing (since $\Sigma$ is a diffeomorphism), \eqref{rel3} implies that
\begin{equation}\label{rel4}
  |\nabla h (x)| \sim  |\nabla \tilde h (\Sigma(x))|, \quad x \in [0,1]^{n-2}
\end{equation}
where the implicit constant depends on $\max_{x \in [0,1]^{n-2}}\{|\nabla \Sigma(x)|\}$ and $\min_{x \in [0,1]^{n-2}}\{|\nabla \Sigma(x)|\}$.
Recalling again from Remark  \ref{basis} that
$y-z$ can be uniquely written as
$$  y-z =  a_0 \nu^\ast + a_1 \tilde e_1 + \cdots + a_{n-2} \tilde e_{n-2} + a_{n-1} \nu
$$
where $a_i\in \R$ for $i=0, \cdots, n-1$.
Moreover, recall that $\nabla_{\tilde e_i}\tilde e_j= \sum_{k=1}^{n-2}\Gamma_{i,j}^k \tilde e_k$ where $\Gamma_{i,j}^k$ are the Christoffel symbols determined by $\Sigma$. Let
\begin{equation}\label{M1}
  M_1:= \max_{w \in \Sigma}\max_{i,j,k=1, \cdots, n-2} |\Gamma_{i,j}^k(w)|.
\end{equation}
%Moreover, by Lemma \ref{dual}, we know $\Sigma^\ast$ has positive principal curvatures with the smallest one at least $\frac{1}{\bar \kappa}$, which implies that the second fundamental form
%\begin{equation}\label{II}
%   T\Sigma^\ast \ni \zeta \mapsto \nabla_\zeta \nu^\ast \cdot \zeta \ge \frac{1}{\bar \kappa}|\zeta|^2.
%\end{equation}
Let
$$  M_2 : = \frac{C_\Sigma}{16n M_1 \max\{1,\bar\kappa^2\}}$$
where $\bar\kappa$ is in \eqref{bdds} and $C_\Sigma$ is from Remark \ref{basis}.
We consider the following two cases separately.

\medskip
\emph{Case 1. $|a_i| \ge M_2|y-z|$ for some $i=1,\cdots, n-1$}. Then recalling \eqref{th},
we deduce that
$$ |h(x)| = |\tilde h (w)|  \ge   |(y-z) \cdot \tilde e_i| = |a_i| \ge M_2|y-z|  \sim |y-z|.$$
Thus \eqref{cine2} holds.

\medskip
\emph{Case 2. $|a_i| < M_2|y-z|$ for all $i=1,\cdots, n-1$}. Noting that $\Sigma_{i=0}^{n-1}a_i^2 \ge C_\Sigma |y-z|^2$, this implies that
\begin{equation}\label{a0}
  |a_0| \ge \frac12|y-z|.
\end{equation}
In this case, we show $|\nabla_\xi h(x)|\gtrsim |y-z|$. It suffices to check $\xi = \partial_i$ for $i=1, \cdots, n-2$, which, combining \eqref{chain} and \eqref{rel4}, further reduces to checking
$$   |\nabla_{\tilde e_i} \tilde h(w)|\gtrsim |y-z|, \quad  i= 1, \cdots,  n-2.   $$
By \eqref{M1}, we compute
\begin{align*}
  |(\nabla_{\tilde e_i}\tilde e_i)\cdot (y-z)| & \le \sum_{k=1}^{n-2} |\Gamma_{ii}^k \tilde e_k \cdot  (y-z)| \le \sum_{k=1}^{n-2} |\Gamma_{ii}^k a_k \tilde e_k| \le (n-2)M_1M_2|y-z|.
\end{align*}
Recalling $\tilde h$ in \eqref{th},
we have
\begin{equation*}
\nabla_{\tilde e_i} \tilde h (w)= (\nabla_{\tilde e_i}[(y-z) \cdot \tilde e_1], \cdots, \nabla_{\tilde e_i}[(y-z) \cdot \tilde e_{n-2}], \nabla_{\tilde e_i}[ (y-z) \cdot\nu]), \quad i=1, \cdots, n-2.
\end{equation*}
We compute
\begin{align*}
  |\nabla_{\tilde e_i} \tilde h(w)| & \ge |\nabla_{\tilde e_i}[(y-z) \cdot \tilde e_i]| =|[\nabla_{\tilde e_i}(y-z)]\cdot \tilde e_i + (\nabla_{\tilde e_i}\tilde e_i)\cdot (y-z)]|  \\
   & \ge  |\nabla_{\tilde e_i}[a_0\nu^\ast+[(y-z)-a_0\nu^\ast]]\cdot \tilde e_i| - |(\nabla_{\tilde e_i}\tilde e_i)\cdot(y-z)| \\
   & \ge |\nabla_{\tilde e_i}(a_0\nu^\ast)\cdot  \tilde e_i| - \sum_{j=1}^{n-2}|(\nabla_{\tilde e_i}a_je_j) \cdot \tilde e_i|- |(\nabla_{\tilde e_i}a_{n-1}\nu) \cdot \tilde e_i| - (n-2)M_1M_2|y-z|\\
   & \ge |a_0||\nabla_{\tilde e_i}\nu^\ast \cdot \tilde e_i| - \sum_{j=1}^{n-2}|a_j||(\nabla_{\tilde e_i}\tilde e_j) \cdot \tilde e_i| -|a_{n-1}||(\nabla_{\tilde e_i}\nu) \cdot \tilde e_i| - (n-2)M_1M_2|y-z|  \\
   & \ge \frac{1}{2 \bar \kappa}|y-z| - (n-2)M_1M_2 |y-z|- \bar \kappa M_2 |y-z| -(n-2)M_1M_2|y-z| \\
   & \ge \frac{1}{4 \bar \kappa}|y-z|
\end{align*}
where in the last two equations we use \eqref{a0}, \eqref{II} and \eqref{IIast}.
Thus \eqref{cine2} is true in this case. Combining two cases, we know (3) in Definition \ref{cine} holds.




\medskip

Finally, we show (2) in Definition \ref{cine}. Indeed, combining \eqref{dou1} and \eqref{cine2}, we deduce that
\begin{equation*}
 |y-z| \lesssim \inf_{x \in [0,1]^{n-2}, \xi \in S^{n-3}} \{ |h(x)| + |\nabla_\xi h(x)| \} \le  \|f_y-f_{z}\|_{C^1([0,1]^{n-1})} \lesssim |y-z|,
\end{equation*}
which implies that (2) in Definition \ref{cine} holds and
$F$ is a doubling space with the same doubling constant as $\mathbb{R}^{n}$. In particular, $F$ is biLipschitz homeomorphic to $Z$.


The proof is complete.
\end{proof}


Finally, Lemma \ref{pairvolume2} holds as a direct consequence of Lemma \ref{fg} and Lemma \ref{cineprop}.




\section{Proof of Theorem \ref{planemain}} \label{s3}
We show Theorem \ref{planemain} in this section. As explained in the introduction, we first prove Theorem \ref{thmconfig} and use Theorem \ref{thmconfig} to obtain Theorem \ref{planemain}.
To this end,
we introduce some definitions.


\begin{definition}[$(\delta,s)$-set] \label{dst}
Let $s \geq 0$, $C > 0$, $\delta> 0$ and $(E,d)$ be a metric space. A bounded set $A \subset E$ is called a \emph{$(\delta,s,C)$-set} if
\begin{displaymath}
|A \cap B(x,r)|_{\delta} \leq Cr^{s}|A|_{\delta}, \qquad x \in E, \, r \geq \delta. \end{displaymath}
  \end{definition}

\begin{remark} \label{deltaset}
  \rm We need following properties of $(\delta,s,C)$-set.

  (i) If $A$ is a non-empty $(\delta,s,C)$-set, then $|A|_{\delta} \geq C^{-1}\delta^{-s}$. This follows by applying the defining condition at scale $r = \delta$.

  (ii) If $A$ is a non-empty $(\delta,s,C)$-set and $A' \subset A$ satisfies $|A'|_{\delta} \ge D^{-1}|A|_{\delta}$ for some $D \ge 1$, then
  \begin{displaymath}
  |A' \cap B(x,r)|_{\delta} \leq |A \cap B(x,r)|_{\delta} \le Cr^{s}|A|_{\delta} \le CDr^{s}|A'|_{\delta}, \qquad x \in E, \, r \geq \delta, \end{displaymath}
  which implies that $A'$ is a $(\delta,s,CD)$-set.
\end{remark}
%The following observations are useful to keep in mind about $(\delta,s,C)$-sets. First, Second, a $(\delta,s,C)$-set is a $(\delta,t,C)$-set for all $0 \leq t \leq s$.

We then introduce the definition of a configuration.
In the following, for all $(f,p) \in C^2([0,1]^{n-2},[0,1]^{n-1}) \times [0,1]^{2n-3}$,
$$\Pi_{1} \colon C^2([0,1]^{n-2},[0,1]^{n-1}) \times [0,1]^{2n-3}  \to C^2([0,1]^{n-2},[0,1]^{n-1}), \quad \Pi_{1}(f,p) := f $$
and
$$\Pi_{2} \colon C^2([0,1]^{n-2},[0,1]^{n-1}) \times [0,1]^{2n-3} \to [0,1]^{2n-3} , \quad \Pi_{2}(f,p) := p $$
stands for the projection to the first variable and to the second variable respectively.
%and $U \subset [0,1]^{n-1}$ is a dyadic cube.



\begin{definition}[Configuration]\label{d:config}
Let $s,t \in (0,n-2]$, $C>0$, and $\delta> 0$. A \emph{$(\delta,s,t,C)$-configuration} is a set $\Omega \subset C^2([0,1]^{n-2},[0,1]^{n-1}) \times \mathbb{R}^{2n-3}$ such that
$$F := \Pi_{1}(\Omega)$$
 is a non-empty $\delta$-separated $(\delta,t,C)$-subset of $C^2([0,1]^{n-2},[0,1]^{n-1})$ consisting of a cinematic family of maps with
cinematic constant $K$ and doubling constant $D$, and for each $f \in F$
$$E(f) := \Pi_{2}((\{f\} \times \mathbb{R}^{2n-3}) \cap \Omega)$$
is the subset of $\mbox{graph}(f)$
such that
\begin{equation}\label{xf}
  E(f) = \{(x,f(x)) \in [0,1]^{2n-3}: x \in X_f\} = (X_f \times [0,1]^{n-1}) \cap \mbox{graph}(f)
\end{equation}
where $X_f$ is a non-empty $(\delta,s,C)$-subset of $[0,1]^{n-2}$.
Also, let
 \begin{displaymath} \label{cale}
 \mathcal{E} := \bigcup_{f \in F} E(f).
 \end{displaymath}
 Additionally, we require that the sets $X_f$ have constant $\delta$-covering number: there exists $M \geq 1$ such that $|X_f|_\delta = M \ge C^{-1}\delta^{-s}$ for all $f \in F$.
%When $U=[0,1]^{n-1}$, we call $\Omega$ a $(\delta,s,t,C)_{n-1}$-configuration for simplicity.
%If the constant $M$ is worth emphasising, we will call $\Omega$ a $(\delta,s,t,C,M)_{n-1}$-configuration.
\end{definition}

We need some auxiliary lemmas.

\begin{lemma} \label{shape}
  Let $F$ be the same cinematic family in Lemma \ref{fg}. For $\epsilon>0$, $0<s \le n-2$ and $X \subset [0,1]^{n-2}$ being a $(\delta,s,\delta^{-\epsilon})$-set with $|X| \le \delta^{-s}$, Denote by $P_{f,g}$ the orthogonal projection of
  $f^\delta([0,1]^{n-2}) \cap g^\delta([0,1]^{n-2})$ to $[0,1]^{n-2}$. Then
  \begin{equation}\label{prfg}
   |X^\delta \cap P_{f,g}|_\delta \lesssim \delta^{-\epsilon} \frac{1}{\|f-g\|_{C^2([0,1]^{n-2})}^s}.
  \end{equation}
\end{lemma}

\begin{proof}
  Let $d=\|f-g\|_{C^2([0,1]^{n-2})}$ and
  $\calU$ be the collection of dyadic cubes $U \subset [0,1]^{n-2}$ from the decomposition of $[0,1]^{n-2}$ in Lemma \ref{fg}. We infer that
  $$  P_{f,g} = \bigcup_{U \in \calU} P_{f,g}^U $$
  where $P_{f,g}^U = P_{f,g} \cap [0,1]^{n-2}$.

  To prove the lemma, it suffices to show
  \begin{equation}\label{covering1}
    |X^\delta \cap P_{f,g}^U|_\delta \lesssim \delta^{-\epsilon} \frac{\delta^{-(n-2)}}{{d}^s}
  \end{equation}
  for each $U \in \calU$.

  Recalling from \eqref{ball} that $P_{f,g}^U \subset  B(x,16K^2\delta/t)$ for some $x \in [0,1]^{n-2}$, we deduce that
  \begin{align*}
  |X^\delta \cap P_{f,g}^U|_\delta & \le  |X^\delta \cap B(x, 16K^2\frac{\delta}{ d })|_\delta \le \delta^{-\epsilon}( 16K^2\frac{\delta}{ d })^{s} |X|  \\
   & \lesssim \delta^{-\epsilon} \frac{\delta^s}{d^s}\delta^{-s} \le \delta^{-\epsilon} \frac{1}{\|f-g\|_{C^2([0,1]^{n-1})}^s},
\end{align*}
which gives \eqref{covering1} as desired.
\end{proof}

Recall that for any set $A \subset \R^l$, $A^\delta$ is the $\delta$-neighbourhood of $A$ in $\R^{l}$. For sets $E(f)$ defined in \eqref{xf}, we write $E^\delta(f)$ instead of $E(f)^\delta$.

\begin{lemma} \label{enbhd} For any $\epsilon>0$ and $\delta >0$, we have
  \begin{equation}\label{3ez}
    \mathcal{L}^{2n-3}(E^\delta(f) \cap E^\delta(g)) \lesssim \delta^{-\epsilon} \frac{\delta^{2n-3}}{\|f-g\|_{C^2([0,1]^{n-2})}^s}.
  \end{equation}
\end{lemma}


\begin{proof}
Write $\|f-g\|_{C^2([0,1]^{n-2})} = d$. Thanks to Lemma \ref{nbhd}, letting $P_{E^\delta(f) \cap E^\delta(g)}$ be the orthogonal projection of $E^\delta(f) \cap E^\delta(g)$ to $[0,1]^{n-2}$,
it suffices to show,
\begin{equation}
    \mathcal{L}^{n-2}(P_{E^\delta(f) \cap E^\delta(g)}) \lesssim \delta^{-\epsilon} \frac{\delta^{n-2}}{d^s}.
  \end{equation}
Noting that $\delta$-balls in $\mathbb{R}^{n-2}$ has measure $\sim \delta^{n-2}$, it further suffices to show,
   $$ |P_{E^\delta(f) \cap E^\delta(g)}|_{\delta} \lesssim \delta^{-\epsilon} \frac{1}{d^s}. $$
   Observing that
   $$P_{E^\delta(f) \cap E^\delta(g)} \subset X_f^\delta \cap P_{f,g} $$
   where $P_{f,g}$ is as in Lemma \ref{shape},
   applying Lemma \ref{shape}, we complete the proof.
\end{proof}

The next theorem provides a lower bound of configurations, which implies Theorem \ref{thmconfig}.


\begin{thm} \label{deltaconfig}
  Let $0<s  \le s' \le n-2$, $\epsilon>0$. Then there exists $\delta_0=\delta_0(\epsilon)>0$ such that the following holds for all $0 < \delta < \delta_0$: for any $(\delta,s',s,\delta^{-\epsilon})$-configuration $\Omega$,
  $$   |\calE|_{\delta} \gtrsim \delta^{7\epsilon-s-s'}$$
  where $\calE$ is in \eqref{cale}.
\end{thm}
\begin{proof}
  Let $C$ be the implicit constant in \eqref{3ez}. The relationship between $\delta_0=\delta_0(\epsilon)$ and $\epsilon$ in this theorem is
  \begin{equation}\label{4ez}
    \log(1/\delta_0) < \delta_0^{-\epsilon}.
  \end{equation}
  Let $\mathcal{E}^\delta$ be the $\delta$-neighbourhood of $\mathcal{E}$. Then it suffices to show $$\mathcal{L}^n(\mathcal{E}^\delta) \gtrsim \delta^{16\epsilon +2n-3- s-s'}.$$
  Recall that
  $$  \mathcal{E}^\delta =  \bigcup_{f \in F} E^\delta(f).$$
  Since we are finding a lower bound, it suffices to assume $|F|$ have an upper bound. Recalling that $F$ is a $\delta$-separated $(\delta,s,\delta^{-\epsilon})$-set, we can assume
  \begin{equation}\label{cardf}
    \delta^{\epsilon} \delta^{-s}\lesssim |F|\lesssim \delta^{\epsilon} \delta^{-s}
  \end{equation}
  in the rest of the proof.
  Using Cauthy-Schwartz type inequality for Lebsegue measure (c.f. \cite[Lemma4.2]{LIU2025}), we have
  \begin{equation}\label{cs}
    \mathcal{L}^{2n-3}(\mathcal{E}^\delta) =  \mathcal{L}^{2n-3}(\bigcup_{f \in F} E^\delta(f)) \ge \frac{[\sum_{f \in F}\mathcal{L}^{2n-3}( E^\delta(f))]^2}{\sum_{f ,g\in F}\mathcal{L}^{2n-3}( E^\delta(f)\cap  E^\delta(g))}.
  \end{equation}
  Recalling $|E^\delta(f)|_{\delta}=M \gtrsim \delta^{\epsilon-s'}$, we have
  $ \mathcal{L}^{2n-3}(E^\delta(f)) \sim M \delta^{2n-3} \gtrsim \delta^{\epsilon} \delta^{2n-3-s'}$ for all $f \in F$. Combining \eqref{cardf}, we deduce that
  $$ [\sum_{f \in F}\mathcal{L}^{2n-3}( E^\delta(f))]^2 \gtrsim [ \delta^{2\epsilon +2n-3 -s-s'} ]^2.  $$
  It suffices to show
  \begin{equation}\label{e2}
    \sum_{f ,g\in F}\mathcal{L}^{2n-3}( E^\delta(f)\cap  E^\delta(g)) \lesssim \delta^{-3\epsilon}\delta^{2n-3-s-s'}
  \end{equation}
  For each $f \in F$ and $i \in  \mathbb{N}$, define
   $$  F_i(f):= \{ g \in F : \|f-g\|_{C^2([0,1]^{n-2})}  \in (2^{-(i+1)},2^{-i}]  \} . $$
   Note that if $i \ge \log (1/\delta)$, then $F_i(f) =\{f\}$ due to $F$ being $\delta$-separated.
   Since $F$ is a $\delta$-separated $(\delta,s,\delta^{-\epsilon})$-set, we have $|F_i(f)| \le 2^{-is}|F|$.
   Using Lemma \ref{enbhd}, we have
   \begin{align}\label{l23}
    \sum_{f \in F}\sum_{g \in F, f\ne g}  \mathcal{L}^{2n-3}( E^\delta(f)\cap  E^\delta(g)) &= \sum_{f \in F}\sum_{i =1}^{\log (1/\delta)}\sum_{g \in F_i(f)}  \mathcal{L}^{2n-3}( E^\delta(f)\cap  E^\delta(g)) \notag\\
    &\overset{\eqref{3ez}}{\lesssim} \delta^{-\epsilon} \sum_{f \in F}\sum_{i=1}^{\log (1/\delta)}2^{-is}|F| \frac{\delta^{2n-3}}{2^{-(i+1)s'}} \notag \\
    & \lesssim  \delta^{-\epsilon}|F|\delta^{2n-3}\sum_{f \in F} \sum_{i=1}^{\log (1/\delta)} 2^{(s'-s)i} \notag\\
    & \lesssim \delta^{-\epsilon}|F|^2 \delta^{2n-3} \log (1/\delta)\delta^{s-s'}
     \le  \delta^{-2\epsilon}\delta^{2n-3-s-s'}
   \end{align}
   where in the second last inequality we note that $s'\ge s$ and $$\sum_{i=1}^{\log (1/\delta)} 2^{(s'-s)i} \le \log (1/\delta) 2^{(s'-s)\log (1/\delta) } = \log (1/\delta) \delta^{s-s'},$$
   and in the last inequality we recall $|F| \le \delta^{-s}$ and \eqref{4ez}.
   Finally, noticing that
   $$  \sum_{f ,g\in F}\mathcal{L}^{2n-3}( E^\delta(f)\cap  E^\delta(g)) = \sum_{f \in F} \mathcal{L}^{2n-3}( E^\delta(f))  + \sum_{f \in F}\sum_{g \in F, g \ne f}  \mathcal{L}^{2n-3}(E^\delta(f)\cap  E^\delta(g))  $$
   and
   $$  \sum_{f \in F} \mathcal{L}^{2n-3}( E^\delta(f)) \lesssim|F|  \delta^{2n-3} M \le \delta^{2n-3-s-s'} \le  \delta^{2n-3-s-s'},  $$
   combining \eqref{l23}, we conclude \eqref{e2} as desired.
\end{proof}


We are ready to show Theorem \ref{planemain}.

\begin{proof}[Proof of Theorem \ref{planemain}]
    Let $Z \subset \R^n$ be analytic with $\dim Z \le n-2$. Since we can always find a sequence of compact subsets $Z_i \subset Z$ such that $\lim_{i \to \infty}\dim Z_i=\dim Z$ and the dimension of the orthogonal projection of a set is invariant under scalings and translations. We could assume $Z \subset B(o,1/2)$.
    Thanks to Proposition \ref{cineprop}, we know that maps $\{f_z\}_{z \in Z}$ defined in \eqref{fw} form a cinematic family of maps.
    %Thus using Theorem \ref{thmconfig} and by a parallel proof of \cite[Section 3.2]{orponen2023hausdorff} or \cite[Proof of Theorem 1.13, Section 5]{LIU2025}, we obtain Theorem \ref{planemain}. We omit the detail.
     We then prove by contradiction. Assume there exists $0<s< t=\dim Z$ and $X_s \subset [0,1]^{n-1}$ such that
  $$ \dim X_s > s $$
  and
  \begin{equation}\label{sdim}
    \dim(\pi_{T_{\Sigma(x)}S^{n-1}}(Z) ) < s, \quad x \in X_s.
  \end{equation}
   We further formulate this statement in a discrete way.
  Fix $s'\in (s, \min\{\dim X_s, t\})$.
  Then the $s'$-dimensional Hausdorff content $H: = \mathcal{H}^{s'}_\infty(Z) >0$
  and
  $$\mathcal{H}^{s'}_\infty( X_s ) >0. $$
  Moreover, by \eqref{sdim}, we know
  $$ \mathcal{H}^{s}(\pi_{T_{\Sigma(x)}S^{n-1}}(Z)) =0, \quad x \in X_s.   $$
  For $\epsilon>0$ and $\delta_0=\delta_0(\epsilon)$, we require Theorem \ref{deltaconfig} holds for all $\delta< \delta_0$. Moreover, by decreasing $\delta_0$, we also require that
  \begin{equation}\label{choiceez}
   \epsilon < \frac{s'-s}{24} \mbox{ and } \delta_0^\epsilon < \min\{ [\log(1/\delta_0)]^{-2}, C_1^{-1}, C_2^{-1}\}
  \end{equation}
  where $C_1,C_2>1$ are the two implicit constant appeared in \eqref{lob1} and \eqref{lob2} respectively.


Fix $t' \in(s',t)$. Thanks to Frostman's lemma, we can find a probability measure $\mu$ with spt$\mu\subset Z$ satisfying $\mu (B(z,r)) <  r^{t'}$ for all $z \in \R^n$ and $r>0$. For each $x \in X_s$, denote by $\mu_x$ the push-forward of $\mu$ under $\pi_{T_{\Sigma(x)}S^{n-1}}$. Thus $\mu_x$ is a probability measure supported on
$\pi_{T_{\Sigma(x)}S^{n-1}}(Z)$, and recalling \eqref{sdim}, we deduce that
$$\calH^{s}({\rm spt}\mu_x)  \le \calH^{s}(\pi_{T_{\Sigma(x)}S^{n-1}}(Z))=0.$$
   Fix $k_0 \in\mathbb{N}$ such that $2^{-k_0} < \delta_0$. For each $x \in X_s$, let $\mathcal{U}_x = \{ B_{x,i}^{n-1}\}_{i \in\mathcal{I}_x}$ be a cover of $\pi_{T_{\Sigma(x)}S^{n-1}}(Z)$($\subset\R^{n-1}=T_{\Sigma(x)}S^{n-1}$)
by $(n-1)$-dimensional balls with dyadic radii $\le 2^{-k_0}$ for all $i \in \mathcal{Q}_x$ such that
$$  \sum_{i \in \mathcal{I}_x}\mu_x(B_{x,i}^{n-1}) =1 \mbox{ and }  \sum_{i \in \mathcal{I}_x} |\diam B_{x,i}^{n-1}|^{s} \le 1.     $$
Write $\mathcal{I}_x^k:= \{ i \in \mathcal{I}_x : | I_x^i|\in [2^{-k},2^{-(k-1)})\}$ for all $k \ge k_0$.
By pigeonhole principle, for each $x \in X_s$, there exists $k_x \ge k_0$ such that
$$    \sum_{i \in \mathcal{I}_x^{k_x}}\mu_{x}(B_{x,i}^{n-1})  \gtrsim k_x^{-2}\mu_{x}({\rm spt}\mu_{x} ) = k_x^{-2} \mbox{ and } |\mathcal{I}_x^{k_x}|  \overset{\eqref{sdim}}{\le} 2^{sk_x}.  $$
By using a second pigeonhole principle, we can find $k_1 \ge k_0$, and a subset $\bar X_s \subset X_s$ such that
  $$  \mathcal{H}^{s'}_\infty(\bar X_s)  \gtrsim k_1^{-2} \mbox{ and } k_x \equiv k_1 \quad x \in \bar X_s.  $$
Let $\delta: = 2^{-k_1}$.
  By \cite[Lemma 2.7]{orponen2023hausdorff}, we know that there exists a $\delta$-separated $(\delta,s',C)$-set $Y_s \subset \bar X_s$ satisfying
  \begin{equation}\label{cardy}
    \delta^{-s'} \lesssim |Y_s| \leq \delta^{-s'}.
  \end{equation}
  As a result, we have
  \begin{equation}\label{upb2}
    \sum_{i \in \mathcal{I}_x^{k_1}}\mu_{x}(B_{x,i}^{n-1}) \gtrsim k_1^{-2} = [\log(1/\delta)]^{-2}  \mbox{ and } |\mathcal{I}_x^{k_1}|  \le \delta^{-s} \quad x \in Y_s.
  \end{equation}
For each $x\in Y_s$, we relate $Z$ to the set $\cup_{i \in \mathcal{I}_x^{k_1}}B_{x,i}^{n-1}$. Let $\calQ(n,\delta)$ be the family of $n$-cubes $Q^n \subset\R^n$ with side length $\delta$.
Define
\begin{equation}\label{qxk}
  \calQ_x : = \{Q^n \in \calQ(n,\delta): \mu(Q^n \cap \pi_{T_{\Sigma(x)}S^{n-1}}^{-1}(\cup_{i \in \mathcal{I}_x^{k_1}}B_{x,i}^{n-1})) >0\},
\end{equation}
and
\begin{equation}\label{zxk}
  \calZ_x : =  \bigcup_{Q^n \in \calQ_x}Q^n \subset \R^n, \quad \forall x \in Y_s.
\end{equation}
We have
$$ \pi_{T_{\Sigma(x)}S^{n-1}}^{-1}(\cup_{i \in \mathcal{I}_x^{k_1}}B_{x,i}^{n-1})) \subset \calZ_x $$
and by the first inequality in \eqref{upb2}, we deduce that
\begin{equation}\label{muzx}
  \mu(\calZ_x) \ge  \mu(\pi_{T_{\Sigma(x)}S^{n-1}}^{-1}(\cup_{i \in \mathcal{I}_x^{k_1}}B_{x,i}^{n-1})) = \mu_x(\cup_{i \in \mathcal{I}_x^{k_1}}B_{x,i}^{n-1}))\gtrsim [\log(1/\delta)]^{-2}, \quad \forall x \in Y_s.
\end{equation}
On the other hand, since each ball $B_{x,i}^{n-1}$ with $i \in \mathcal{I}_x^{k_1}$ has radius $2^{-k_1}=\delta$ and $\calZ_x$ is a union of $\delta$-cubes, we know
\begin{equation}\label{inclzx}
  \pi_{T_{\Sigma(x)}S^{n-1}}(\calZ_x) \subset \bigcup_{i \in \mathcal{I}_x^{k_1}} 2B_{x,i}^{n-1}
\end{equation}
where $2B_{x,i}^{n-1}$ is the ball with same center and twice the radius as $B_{x,i}^{n-1}$.

Let
$$ \calQ:= \bigcup_{x \in Y_s}\calQ_x \mbox{ and } \calZ:= \bigcup_{x \in Y_s}\calZ_x.   $$
We have
\begin{align*}
  \sum_{Q^n \in \calQ}\mu(Q^n)|\{x & \in Y_s: Q^n \in \calQ_x\}| =  \sum_{x \in Y_s} \mu(\calZ_x)\\
  &\overset{\eqref{cardy}}{\gtrsim} \delta^{-s'} \min_{x \in Y_s}\{\mu(\calZ_x)\}
    \overset{\eqref{muzx}}{\gtrsim} \delta^{-s'}[\log(1/\delta)]^{-2}\overset{\eqref{choiceez}}{\gtrsim} \delta^{\epsilon}\delta^{-s'}.
\end{align*}
Since $\{x  \in Y_s: Q^n \in \calQ_x\} \subset Y_s$, we have
$$  |\{x  \in Y_s: Q^n \in \calQ_x\}| \le |Y_s| \overset{\eqref{cardy}} \le \delta^{-s'}, \quad \forall x \in Y_s.  $$
Combining the above two estimates, we infer that there exists a subfamily $\widetilde{\calQ} \subset \calQ$ with
$$ \mu(\widetilde{\calZ}) \gtrsim \delta^{2\epsilon} \mbox{ where $\widetilde{\calZ}:= \bigcup_{Q^n \in \widetilde{\calQ}}Q^n \subset \R^n$}$$
such that each $Q^n \in \widetilde{\calQ}$ satisfies
\begin{equation}\label{cardysz}
   |\{x  \in Y_s: Q^n \in \calQ_x\}| \gtrsim \delta^{2\epsilon}\delta^{-s'}.
\end{equation}
Recalling that spt$\mu=Z$, we have $\mu(\widetilde{\calZ} \cap Z)= \mu(\widetilde{\calZ}) \gtrsim \delta^{2\epsilon}$. Since $\mu$ is a $t'$-dimensional Frostman measure and $s' < t'$, we deduce that
$$\calH_\infty^{s'}(\widetilde{\calZ} \cap Z) \gtrsim \calH_\infty^{t'}(\widetilde{\calZ} \cap Z) \gtrsim  \delta^{2\epsilon}.$$
 Thanks to \cite[Lemma 2.7]{orponen2023hausdorff}, we can find a $\delta$-separated $(\delta,s',\delta^{-2\epsilon})$-set $\widetilde{Z} \subset \widetilde{\calZ} \cap Z$ with $|\widetilde{Z}| \gtrsim \delta^{2\epsilon} \delta^{-s'}$.
 For each $z \in \widetilde{Z}$, there exists a $Q^n_z \subset \widetilde{Q}$ such that $z \in Q^n_z$. Define
 $$ Y_{s}(z):= \{x  \in Y_s: Q^n_z \in \calQ_x\} \subset Y_s, \quad \forall z\in \widetilde{Z}. $$
 By \eqref{cardysz}, we have
 $$ |Y_{s}(z)|  \gtrsim  \delta^{2\epsilon}\delta^{-s'} \gtrsim \delta^{2\epsilon}|Y|. $$
 In particular, recalling Remark \ref{deltaset} (ii), $Y_{s}(z)$ a is $(\delta,s',\delta^{-3\epsilon})$-set.


 To apply Theorem \ref{deltaconfig}, we construct a $(\delta,s',s',\delta^{-3\epsilon})$-configuration.
 For $z \in \widetilde{Z}$, let
 $$   E(f_z):=\bigcup_{x \in Y_{s}(z)} (x, f_z(x))  \subset [0,1]^{2n-3}.$$
 and
 $$ \Omega:= \bigcup_{z \in \widetilde{Z}} (f_z,E(f_z)).$$
 Then $\Omega$ is a $(\delta,s',s',\delta^{-3\epsilon})$-configuration.
 Recall $\mathcal{E} := \bigcup_{z \in \widetilde{Z}} E(f_z)$. Applying Theorem \ref{deltaconfig}, we know
 \begin{equation}\label{lob1}
   |\mathcal{E}|_\delta  \ge C_1^{-1}\delta^{21\epsilon-2s'}
 \end{equation}
 for some $C_1>1$.

On the other hand,
noting that $x\in Y_s(z)$ implies $z\in Q_z^n \subset \calZ_x$, we have
$$ \mathcal{E} = \bigcup_{z \in \widetilde{Z}} \bigcup_{x \in Y_s(z)}(x,f_{z}(x)) \subset \bigcup_{x \in Y} \bigcup_{z \in \calZ_x} (x,f_{z}(x)) = \bigcup_{x \in Y} \bigcup_{z \in \calZ_x} (x,\pi_{T_{\Sigma(x)}S^{n-1}}(z)) = \bigcup_{x \in Y}\pi_{T_{\Sigma(x)}S^{n-1}}(\calZ_x).$$
Recalling \eqref{upb2} and \eqref{inclzx},
we deduce
\begin{equation}\label{lob2}
  |\mathcal{E}|_\delta \le \sum_{x \in Y} |\pi_{T_{\Sigma(x)}S^{n-1}}(\calZ_x)|_\delta \le C_2|Y|\delta^{-s}\le C_2\delta^{-s-s'} .
\end{equation}
This contradicts to \eqref{lob1}, since $\epsilon < 24^{-1}(s'-s)$ by our choice of $\epsilon$ in \eqref{choiceez}. The proof is complete.
\end{proof}




\section{Line-cone Incidence Estimate}\label{s4}

We prove Theorem \ref{lineconen} in this section and we recall the statement of Theorem \ref{lineconen} here.

\medskip
\noindent
{\bf Theorem 1.28.}
Let $\Sigma$ be as in Theorem \ref{cormain}.
If a line $L \subset \R^n$ satisfies $\min_{p \in \calC_z}\{\angle(L,T_p\calC_z)\}  \ge a$ for some $a \gg \delta >0$, then
  $$\calL^{n}(\calC_z^\delta \cap L^\delta) \lesssim_a \delta^n $$
  where we recall the cone $\calC_z$ with apex $z$ and cross-section $\Sigma$ is the set
$$\calC_z = \bigcup_{p \in \Sigma} L_{z,p} = \bigcup_{p \in \Sigma, r \in [-1,1]} z + rp \subset \R^n.$$


Before presenting the proof, we make several reductions. First, we can partition $\Sigma$ into small pieces $\Sigma_i$ and show Theorem \ref{lineconen} holds for each cone with cross-section $\Sigma_i$ sufficiently small. Thus in the following of this section,
by partitioning $\Sigma$ into small pieces $\Sigma_i$, we can assume each $\Sigma_i: [0,1]^{n-2} \to S^{n-1}$ has diameter $<\frac{\pi}{100}$. In particular, for any $x,y \in \Sigma_i$, $\angle(x,y) < \frac{\pi}{100}$. Moreover, we require that the smallness of $\Sigma_i$ also guarantees the validity of Lemma \ref{smalla} below.
In the following, we focus on one piece $\Sigma_i$ and write $\Sigma=\Sigma_i$.

Second, noting that
$$ \calL^{n}(\calC_z^\delta \cap L^\delta) = \calL^{n}[(\calC_z^\delta-\{z\}) \cap (L^\delta-\{z\})] = \calL^{n}[\calC_o^\delta \cap (L^\delta-\{z\})], $$
it suffices to show Theorem \ref{lineconen} holds for the cone $\calC_o$ with apex as the origin $o$. In the following, we write $\calC_o = \calC$.
%recalling the cone with apex as the origin $o$ and link $\Sigma$:
%$$\calC_o = \bigcup_{p \in \Sigma} L_{o,p} = \bigcup_{p \in \Sigma, r \in [-1,1]} o + rp.$$

Third, noting that we can decompose $\calC$ as $\calC_+ \cup \calC_-$ where
$$ \calC_\pm:=  \bigcup_{p \in \Sigma, r \in [0,1]}  \pm rp, $$
to show Theorem \ref{lineconen}, it suffices to estimate $\calL^{n}(\calC^\delta_\pm \cap L^\delta)$ separately. Since the proof will be the same for both $\calC_\pm$, we will only focus on $\calC_+$ and write $\calC_+=\calC$ and $L_{o,p}=\cup_{r \in [0,1]}rp$ throughout the rest of this section.
Moreover,
since $\calL^{n}(\calC^\delta \cap L^\delta)$ is invariant under any isometry of $\R^n$, we can always rotate $\calC^\delta$ with respect to $o$ such that one generatrix of $\calC$ is the direction of the first axis $e_1$ and the tangent space of this generatrix is span$\{e_1, \cdots, e_{n-1}\}$. The choice of this generatrix will be in the beginning of the proof of Theorem \ref{lineconen} below \eqref{x} at the end of this section.
By doing this and up to a flipping with respect to span$\{e_1, \cdots, e_{n-1}\}$, the cone $\calC$ can be seen as a graph of a function $h$ from a domain $U \to \R_+$ where $U \subset$ span$\{e_1, \cdots, e_{n-1}\}$ is the orthogonal projection of $\calC$ to span$\{e_1, \cdots, e_{n-1}\}$ (noting that by requiring the diameter of $\Sigma$ sufficiently small we could assume that the projection acting on $\calC$ is a homeomorphism and the function is well-defined on $U$). And in the following, we will consider $\calC^\delta$ as the vertical $\delta$-neighbourhood of $\calC=$ graph$(h)$, i.e.
$$ \calC^\delta = \{ (x,y)\in U \times \R: |y-h(x)|< \delta \}. $$







We first prove the following result.

\begin{lemma} \label{twopt}
Let $\Sigma$ be as in Theorem \ref{planemain}. For  any line $L \subset \R^n$ that is not a generatrix of $\calC$,
$L \cap \calC$ consists of at most $2$ points.
%Moreover, if $L$ is tangent to $\calC_w$, then $L \cap \calC_w$ consists of exactly one point.
\end{lemma}

\begin{proof}
%   Without loss of generality, we can assume $w=o$ is the origin and we write $\calC_o=\calC$ for simplicity.
   Assume $L \cap \calC$ consists of at least three points $x,y,z$ and
   \begin{equation}\label{y}
     y = \eta x+ (1- \eta)z, \mbox{ for some $\eta \in (0,1)$}.
   \end{equation}
    Consider the radial projection $\bf{r}$ at $o$, i.e.
    $${\bf r}(p) = \frac{p}{|p|} \in S^{n-1}, \quad \forall p \in \R^n\setminus \{o\}.$$
    Since $L$ is not a a generatrix of $\calC$, $L \cap \{o\} = \emptyset$. Thus there exists a unique plane $P_L \subset R^n$ passing through $o$ containing $L$.
    One observes that
    $$  {\bf r}(L) = P_L \cap S^{n-1} := S^1_L \mbox{ is a great circle}$$
    and
    $$  {\bf r}(\calC) = {\bf r}(\bigcup_{p \in \Sigma} L_{o,p}) = \bigcup_{p \in \Sigma}L_{o,p} \cap S^{n-1} = \Sigma.$$
    Thus
    \begin{equation}\label{y'}
      \{x',y',z'\} := {\bf r}(\{x,y,z\}) \subset {\bf r}(L \cap \calC) \subset S^1_L \cap \Sigma.
    \end{equation}
    Since $\diam \Sigma< \frac{\pi}{100}$, $|x'-z'|< \frac{\pi}{100}$. Thus the arc $\widehat{x'z'} \subset S^1_L$ with length $=\angle(x',z')< \pi$ is the unique length-minimising geodesic joining $x'$ and $z'$.
    By the condition $\Sigma$ has non-vanishing geodesic curvature ($n=3$) and sectional curvature $>1$ ($n \ge 4$), we know that
    $\Sigma$ is a strictly convex set in $S^{n-1}$ under the standard spherical metric. Thus for any two distinct points on $\Sigma$, the unique length-minimising geodesic segment joining them does not intersect $\Sigma$ except the endpoints.

However, \eqref{y} implies
$$ y' \subset \widehat{x'z'}, $$
and \eqref{y'} implies
$$ \{x',y',z'\}  \subset \widehat{x'z'} \cap \Sigma, $$
which is a contradiction.

%\medskip
%If $L$ is tangent to $\calC$ at some point $q$. Then $L \subset T_q\calC$. Assume $q \in L_{o,p}$
% for some $p \in \Sigma$. We know $L_{o,p} \subset T_q\calC$. Since $L$ is not a generatrix, we know $L \ne L_{o,p}$ and by $L \cap L_{o,p} = q$, we know
% \begin{equation}\label{q}
%   q = L \cap L_{o,p}.
% \end{equation}
% By $L_{o,p} \subset T_q\calC$ and \eqref{q}, one has
% $$ (L_{o,p} \cap \calC) \setminus \{q\} \subset  (T_q\calC \cap \calC) \setminus \{L_{o,p}\}.  $$
% It suffices to check
% $$(T_q\calC \cap \calC) \setminus \{L_{o,p}\} = \emptyset.$$
%  Note that we can write $\calC$ as
%  $$   \calC \setminus \{o\}= \bigcup_{r\in[-1,1]} {\bf r}^{-1}(\Sigma) \cap S^{n-1}_r =: \bigcup_{r\in[-1,1]} \Sigma_r$$
%  where ${\bf r}^{-1}$ is the preimage set of the radial projection and $S^{n-1}_r$ is the sphere centred at $o$ with radius $r$.
%  It reduces to checking that
%  $$(T_q\calC \cap \calC \cap S^{n-1}_r) \setminus \{L_{o,p}\} = (T_q\calC \cap \Sigma_r) \setminus \{rp\} =\emptyset, \quad \forall r \in [-1,1].$$
%
%  On the other hand, noting that for any $q' \in L_{o,p}$, we have $q' \in \Sigma_{|q'|}$ and
%  $$ T_q\calC = T_{q'}\calC \mbox { and } T_{q'}\calC = T_{q'}\Sigma_{|q'|} \oplus \mbox {span}\{q'\} = T_{q'}\Sigma_{|q'|} \oplus \mbox {span}\{p\},  $$
%  combining that $L_{o,p} \subset \mbox {span}\{p\}$, it further reduces to checking that
%  $$  T_{q'}\Sigma_{|q'|} \cap \Sigma_{|q'|} \setminus \{q'\} =\emptyset , \quad \forall q \in L_{o,p}.  $$
%  Note that $\Sigma_r \subset S^{n-1}_r$ is the image of $\Sigma$ under the scaling map $x \mapsto rx$, $x \in \R^n$, which implies that
%  $\Sigma_r$ has sectional curvature $> r^2$. Thus for any $q' \in L_{o,p}$, $\Sigma_{|q'|}$ is strictly convex and the tangent space  $T_{q'}\Sigma_{|q'|}$ of $\Sigma_{|q'|}$ at $q'$ only intersects $\Sigma_{|q'|}$ at $q'$.

  The proof is complete.
\end{proof}

We have the following consequence.

\begin{cor} \label{int2}
   Let $\Sigma$ be as in Theorem \ref{planemain}. Then $L \cap \calC^\delta$ has at most $2$ connected components.
\end{cor}

%\begin{proof}
%    Noting that $\calC^\delta$ is an open set in $\R^n$ and $L \subset \R^n$ is a line, $L \cap \calC^\delta$ consists of open intervals of $L$.
%
%   Assume on the contrary that $L \cap \calC^\delta$ consists at least three disjoint  intervals $I_1,I_2,I_3 \subset L$. We can assume their closures are also disjoint, otherwise for all $\delta'$ slightly smaller than $\delta$, these intervals will have disjoint closures and we derive the contradiction for $\delta'$.
%
%   %Since $\diam \Sigma< \frac{\pi}{100}$ and one generatrix of $\calC$ is the direction $e_1$, we know each generatrix $L_{o,p}$ satisfying $\angle (L_{o,p}, e_n) \in [\frac{\pi}{4}-\frac{\pi}{100}, \frac{\pi}{4}+\frac{\pi}{100}]$
%%   Up to a rotation in $\R^n$, we can assume that all generatrices of $\calC$ is not contained in
%%   the last coordinate span$\{e_n\}$ (recalling, we can arrange ). Thus $\calC$ can be viewed as a graph of the a function in $\R^n$ and we can assume $\calC^\delta$ to be the ``vertical neighbourhood'':
%Note that the vertical neighbourhood $\calC^\delta$  can be written as
%   $$  \calC^\delta= \bigcup_{|\epsilon| < \delta} \calC + \epsilon e_n = \bigcup_{|\epsilon| < \delta} \calC_{\epsilon e_n}. $$
% We have
%   $$  \partial \calC^\delta =\calC_{\delta e_n} \cup \calC_{-\delta e_n} \cup [(\partial U \times \R) \cap \partial \calC^\delta]. $$
%   Let $\partial I_i =\{a_i,b_i\} \subset L$, $i=1,2,3$. We consider the following cases separately.
%
%   \medskip
%   \emph{Case 1. $\{a_i,b_i\}  \subset  \calC_{\delta e_n} \cup \calC_{-\delta e_n}, \ i=1,2,3.$}
%   This implies that either $\calC_{\delta e_n}$ or $\calC_{-\delta e_n}$ would contain more than two points among
%   $\{a_1,b_1, a_2,b_2,a_3,b_3\} \subset L$. Thus $L$ would intersect either $\calC_{\delta e_n}$ or $\calC_{-\delta e_n}$ more than $2$ points, which contradicts to Lemma \ref{twopt}.
%
%\medskip
%   \emph{Case 2. Case 1 does not hold, i.e. $L \cap [(\partial U \times \R) \cap \partial \calC^\delta] \ne \emptyset$.}
%   Let . Then we know $L_\pm \cap [(\partial U \times \R) \cap \partial \calC^\delta] = \emptyset$.
%   As a result,
%   $$ \partial(L_\pm \cap \calC^\delta) \subset \calC_{\delta e_n} \cup \calC_{-\delta e_n}. $$
%   By Case 1, we know $L_\pm \cap \calC^\delta$ has at most two connected components.
%\end{proof}

%The next lemma focuses on the geometric property of the cone $\calC$. Recall that
%$$\kappa_1(y),\cdots, \kappa_{n-2}(y)$$
% are the principal curvatures of $\Sigma \subset S^{n-1}$ at $y \in \Sigma$.
%
%
%\begin{remark} \label{bar}
%  \rm By considering $\Sigma \subset \R^n$ as a codimension $2$ submanifold, we can also compute the principal curvatures $\bar \kappa_i$, $i=1, \cdots,n-2$
%  of $\Sigma$ with respect to the unit normal $\nu$ in $R^n$. We remark that $\bar \kappa_i = \kappa_i$ for all $i=1, \cdots,n-2$.
%  This is because principal curvatures are eigenvalues of the second fundamental form.
%  Thus $\{\kappa_i\}_{i=1, \cdots,n-2}$ are the eigenvalues of the following bilinear form
%  \begin{equation}\label{nabla}
%    T\Sigma\times T\Sigma \ni (X,Y) \mapsto \nabla_X Y \cdot \nu
%  \end{equation}
%  where $\nabla$ is a Levi-Civita connection on the standard sphere $S^{n-1}$
%  and $\{\bar \kappa_i\}_{i=1, \cdots,n-2}$ are the eigenvalues of the following bilinear form
%  \begin{equation}\label{nabla2}
%    T\Sigma\times T\Sigma \ni (X,Y) \mapsto \bar\nabla_X Y\cdot \nu
%  \end{equation}
%  where $\nabla$ is a Levi-Civita connection in $\R^n$.
%
%  Since at each $y\in \Sigma$, we have the relation
%  $$ \bar\nabla_X Y  =  \nabla_X Y + (\nabla_X Y \cdot y)y, $$
%  by noting $y \perp \nu$, we know the forms \eqref{nabla} and \eqref{nabla2} coincide and $\bar \kappa_i = \kappa_i$ for all $i=1, \cdots,n-2$.
%\end{remark}

%\begin{lemma} \label{similar}
%   For any point $y \in \calC \setminus \{o\}$, the principal curvatures of $\calC$ at $y$ are
%   $$ 0, \frac{1}{|y|}\kappa_1(y),\cdots, \frac{1}{|y|}\kappa_{n-2}(y)$$
%   and the principal direction of the value $0$ is the direction of the generatrix passing through $y$, i.e. $\frac{y}{|y|}$.
%\end{lemma}
%
%
%\begin{proof}
%%   Since $\calC_w=\calC_o + w$ for all $w \in \R^n$, and translations in $\R^n$ are isometries which keep the principal curvatures, it suffices to show the case $w=o$ and we write $\calC=\calC_o$.
%
%   Note that for any $y \in \calC$, the generatrix $L_{o,\frac{y}{|y|}}$ which is a line segment is contained in $\calC$, thus the principal curvature with respect to this direction is $0$.
%
%     To see the other principal curvatures, note that we can write $\calC$ as
%     \begin{equation}\label{sigmar}
%       \calC \setminus \{o\}= \bigcup_{r\in[-1,1]} {\bf r}^{-1}(\Sigma) \cap S^{n-1}_r =: \bigcup_{r\in[-1,1]} \Sigma_r
%     \end{equation}
%  where ${\bf r}^{-1}$ is the preimage set of the radial projection and $S^{n-1}_r$ is the sphere centred at $o$ with radius $r$.
%  Thus $\Sigma_r \subset S^{n-1}_r$ is the image of $\Sigma$ under the scaling map $x \mapsto rx$, $x \in \R^n$, which implies the principal curvatures at $y \in \Sigma_{|y|}$ is $\frac{1}{|y|}$ times the principal curvatures at $\frac{y}{|y|}\in\Sigma_{1}=\Sigma$.
%
%  It suffices to show the lemma for $y \in \Sigma \subset \calC$.
%  Recall that at $y  \in \Sigma$, we have the relation
%  $$  \R^n= \mbox{span}\{y,\nu\} \oplus T_y\Sigma$$
%  and $T_y\calC= \mbox{span}\{y\} \oplus T_y\Sigma$ where $\nu$ is the unit normal of $\Sigma \subset S^{n-1}$. Thus $\nu$ is also the unit normal of $\calC \subset \R^n$. Since we already know $\mbox{span}\{y\}$ is one of the  principal direction with principal curvature $0$, it suffices to find principal curvatures with
%   the  principal directions in $T_y\Sigma \subset T_y\calC$, which are
%the eigenvalues of the second fundamental form $\nu$ restricted to $T_y\Sigma$. Using Remark \ref{bar}, we conclude the desired result.
%\end{proof}
%
%\begin{remark} \label{graph}
%   We recall a fact that any manifold $M$ can be a seen as a graph of a function $h$ from a point $x \in M$ locally. And moreover, the eigenvalues of the Hessian $\nabla^2 h(x)$ at $x$ coincides with eigenvalues of the second fundamental form of $M$ at $x$ which are precisely the principal curvatures of $M$ at $x$.
%\end{remark}
%
%
%
%\begin{lemma} \label{small}
%  Let $h: U \to \R$ be the function such that graph$(h)=\calC$.
%  If $\Sigma$ is a sufficiently small piece, then the Hessian of $h$ at $x \in U\setminus\{o\}$ has one eigenvalue $0$ and all other eigenvalues $> \frac{\underline \kappa}{2|x|}$.
%
%  Moreover, for any $a>0$, we can also require the smallness of $\Sigma$ depending on $a$ such that
%  \begin{equation}\label{10a}
%    |\nabla h(x)| \le \frac{a}{10}, \quad \forall x \in U \setminus\{o\}.
%  \end{equation}
%\end{lemma}
%
%\begin{proof}
%  By the convention at the beginning of this section, there is a generatrix $L_{o,p}$ of $\calC$ with direction $e_1$ where $p \in \Sigma$. In particular, $L_{o,p} \subset U$. We first show the lemma for points in this generatrix.
%
%  Let $x \in L_{o,p}$. Then by our convention at the beginning of this section, we know $T_x\calC =$ span$\{e_1,\cdots,e_{n-1}\}$. Thus at $x$, $\calC$ is the graph of the function $h$ from its tangent space. By Remark \ref{graph} and Lemma \ref{similar}, we know at $x \in \calC$, the Hessian of $h$ has one eigenvalue $0$ and all other eigenvalues $\ge \frac{\underline \kappa}{|x|}$. Thus the Lemma holds for $x \in L_{o,p}$.
%
%  We then prove the lemma for other points. To this end, consider the following decomposition of $U$. Recalling from \eqref{sigmar} that
%  $$ \calC = \bigcup_{r \in [-1,1]} \Sigma_r.$$
%  For each $r \in [-1,1]$, let $U_r$ be the orthogonal projection of $\Sigma_r$ to span$\{e_1,\cdots,e_{n-1}\}$. Then, we know
%  $$  U = \bigcup_{r \in [-1,1]} U_r.$$
 %Let $y \in U_1$.
%  Since $\Sigma$ is $C^2$, we know the eigenvalues are continuously depending on $z \in U\setminus\{o\}$, thus if $\Sigma$ is chosen sufficiently small, we know $U_1$ is also small such that all points $y \in U_1$ are sufficiently close to $U_1 \cap L_{o,p}$.
%  As a result, the Hessian of $h$ has one eigenvalue $0$ and all other eigenvalues $\ge \frac{\underline \kappa}{2}$ for all $y \in U_1$.
%
%  For other points $z \in U_r$, since $\frac{z}{|z|} \in U_1$, by Lemma \ref{similar}, we deduce the Hessian of $h$ has one eigenvalue $0$ and all other eigenvalues $\ge \frac{\underline \kappa}{2|z|}$ for all $z \in U_r$.
%
%
%  The proof is complete.
%\end{proof}


\begin{lemma} \label{smalla}
  Let $h: U \to \R$ be the function such that graph$(h)=\calC$. For any $a>0$,
  if $\Sigma$ is a sufficiently small piece depending on $a$, then
  \begin{equation}\label{10a}
    |\nabla h(x)| \le \frac{a}{10}, \quad \forall x \in U \setminus\{o\}.
  \end{equation}
\end{lemma}

\begin{proof}
    By the convention at the beginning of this section, there is a generatrix $L_{o,p}$ of $\calC$ with direction $e_1$ where $p \in \Sigma$. In particular, $L_{o,p} \subset U$. We first show the lemma for points in this generatrix.

  Let $x \in L_{o,p}$. Then by our convention at the beginning of this section, we know $T_x\calC =$ span$\{e_1,\cdots,e_{n-1}\}$. Thus at $x$, graph$(h)=\calC$ is tangent to span$\{e_1,\cdots,e_{n-1}\}$. This implies that $\nabla h(x)=0$. Thus the lemma holds for $x \in L_{o,p}\setminus\{o\}$.

  We then prove the lemma for other points. To this end, consider the following decomposition of $U$. note that we can write $\calC$ as
     \begin{equation}\label{sigmar}
       \calC \setminus \{o\}= \bigcup_{r\in[-1,1]} {\bf r}^{-1}(\Sigma) \cap S^{n-1}_r =: \bigcup_{r\in[-1,1]} \Sigma_r
     \end{equation}
  where ${\bf r}^{-1}$ is the preimage set of the radial projection and $S^{n-1}_r$ is the sphere centred at $o$ with radius $r$.
  For each $r \in [-1,1]$, let $U_r$ be the orthogonal projection of $\Sigma_r$ to span$\{e_1,\cdots,e_{n-1}\}$. Then, we know
  $$  U = \bigcup_{r \in [-1,1]} U_r.$$
    Let $y \in U_r$ for some $r \in [-1,1]$.
  Since $\Sigma$ is $C^2$, we know $\nabla h$ is continuously depending on $U\setminus\{o\}$, thus if $\Sigma$ is chosen sufficiently small, we know $U_r$ is also small such that all points $y \in U_r$ are sufficiently close to $U_r \cap L_{o,p}$.
  As a result, $|\nabla h(y)| \le \frac{a}{10}$ for all $y \in U_r$.

  The proof is complete.
\end{proof}



We are able to show Theorem \ref{lineconen}.

\begin{proof}[Proof of Theorem \ref{lineconen}]
    %It suffices to consider the case $w=o$ and we write $\calC=\calC_o$.
    Recall
    $$ a:= \min_{p \in \Sigma}\{ \angle(L,p) \}. $$
    Note that $L^\delta \cap \calC^\delta$ is the $\delta$-neighbourhood of $L \cap \calC^\delta$.
    By
     Corollary \ref{int2},
     It suffices to show each interval $I \subset L$ in the intersection $L \cap \calC^\delta$ has length upper bound $\lesssim \frac{\delta}{a}$.
     %if $L \cap \calC^\delta$ consists of at most two intervals in $L$, we obtain the desired result. If on the other hand, $L_\pm \cap \calC^\delta$ consists of at most two intervals (recalling $L_\pm$ from Corollary \ref{int2}), noting that
%     $$L^\delta \cap \calC^\delta \subset (L_+^{4\delta} \cap \calC^{4\delta}) \cup (L_-^{4\delta} \cap \calC^{4\delta}),$$
%     we can also obtain that
%     $$ \calL^n(L^\delta \cap \calC^\delta) \le \calL^n(L_+^{4\delta} \cap \calC^{4\delta})+\calL^n(L_-^{4\delta} \cap \calC^{4\delta}) \lesssim \frac{\delta^n}{a}$$
%  as desired.
     Thus,
     we concentrate on a fixed interval $I \subset L$.

    We consider the following two cases depending on the Euclidean distance $d(o,I)$ from $o$ to $I$.

    \medskip
    \emph{Case 1. $I \cap B(o,10 \delta) = \emptyset$.}

    %We choose a maximal $\delta$-separated $(\delta,n-2)$-set $\Sigma_\delta \subset\Sigma$ with $|\Sigma_\delta| \sim \delta^{-(n-2)}$. We observe that
%    $$  \calC^\delta = \bigcup_{x \in \Sigma_\delta} L_{o,x}^\delta. $$
    Since $I \subset L \cap \calC^\delta$, we can find $x\in I \subset L \cap \calC^\delta$ and $y \in \calC$ such that
    $$  x = y + \eta \nu_y  $$
    where $\nu_y$ is the unit normal of $T_y\calC$ and $\eta \in (-\delta, \delta)$.
    Equivalently, we have
    \begin{equation}\label{x}
       |x-y| < \delta, \mbox{ and } x-y \perp T_y\calC.
    \end{equation}
    In particular, since $x \notin B(o,10 \delta)$, $|y| > 9 \delta$ and hence $y \ne o$.

    As mentioned in the beginning of this section, we choose $y$ as the first coordinate and $y = |y|e_1$ and $T_y\calC=$ span$\{e_1, \cdots, e_{n-1}\}$. Thus $\nu_y= e_n$ and $x = y+ x \cdot e_n$.
    %By $\diam \Sigma < \frac{\pi}{100}$, we know all generatrices of $\calC$ is not contained in the last coordinate span$\{e_n\}$.
    And we view $\calC$ as the graph of a function $h:U \to \R$ where we recall $U \subset [0,1]^{n-1}\subset$ span$\{e_1,\cdots,e_{n-1}\}$ is the orthogonal projection of $\calC$ to span$\{e_1,\cdots,e_{n-1}\}$.
    Let $y_\ast:= \frac{y}{|y|} \in \Sigma$. We note that the generatrix $L_{o,y_\ast} \subset$ span$\{e_1\}\cap \calC \subset U$. In particular, $y \in U$. The following figure shows the case when $n=3$ and $L$ passing through $\calC$ (where one has $x=y$).

    \begin{figure}[h!]
\begin{center}
\begin{overpic}[scale = 0.5]{lineconeA}
\end{overpic}
\caption{Cone $\calC$, line $L$, direction $e_1^\ast$ and plane $P_L$}\label{fig1}
\end{center}
\end{figure}


\medskip
    \emph{Subcase 1-1. $L \perp T_y\calC$.}

    Namely, $L$ is parallel to $e_n$ (i.e. $\angle (L,T_y\calC) =\frac{\pi}2$). In this subcase, it is easy to see $|L \cap \calC^\delta|= 2\delta$ since the vertical thickness of $\calC^\delta$ is $2\delta$.


\medskip
    \emph{Subcase 1-2. $\angle (L,T_y\calC) \in [a, \frac{\pi}{2})$.} We refer to Figure \ref{fig2} for an illustration.

        \begin{figure}[h!]
\begin{center}
\begin{overpic}[scale = 0.6]{lineconeB}
\end{overpic}
\caption{Graph$(h_\ast)$, intervals $I$, $I_1$ and $I_2$}\label{fig2}
\end{center}
\end{figure}



     Let $e_1^\ast$ be the direction of the orthogonal projection of $L$ into $T_y\calC=$ span$\{e_1, \cdots, e_{n-1}\}$. Since $x \in L$, we know
     $$L \subset P_L:=\mbox{span}\{e_1^\ast,e_n\}+\{x\}.$$
     By $x-y \perp T_y\calC$ (recalling \eqref{x}),
     we know $x-y$ is parallel to $e_n \in P_L$, thus
     $y \in P_L$.
     In the following, we will focus on the plane $P_L$ and it is convenience to let $y$ be the origin of $P_L$, $e_1^\ast$ be the first coordinate and $e_n$ be the second (vertical) coordinate. Thus $y=(0,0) \in P_L$ and $x$ has coordinate
     $(0,x-y)\in P_L$.

     Define
     $$U_\ast := U \cap [\mbox{span}\{e_1^\ast\}+ \{x\}].$$
     Then $L \cap (U_\ast \times \mbox{span}\{e_n\})$ can be seen as a graph in $P_L$ of a linear function $g: U_\ast \to \R$. Let $h_\ast:= h|_{U_\ast}$.
     We write $s$ as the elements of $U_\ast$ when considering $U_\ast$ as a subset of $\mbox{span}\{e_1^\ast\}+ \{x\} \subset P_L$ and write $se_{1}^\ast$ as the elements of $U_\ast$ when considering $U_\ast$ as a subset of $\R^n$. Note that
     $s=0$ is the point $y$.

     To estimate the length of the interval $I \subset L$,
     let $I_1$ be the orthogonal projection of $I$ to $\mbox{span}\{e_1^\ast\}+ \{x\}$ and $I_2$ be the orthogonal projection of $I$ to $\mbox{span}\{e_n\}$. We have the relation
     \begin{equation}\label{I12}
        |I| = \frac{|I_1|}{\cos \angle(I,I_1)} = \frac{|I_1|}{\cos \angle(L,T_y\calC )}, \quad
     |I| = \frac{|I_2|}{\cos \angle(I,I_2)} = \frac{|I_2|}{\sin \angle(L,T_y\calC)}
     \end{equation}
     and
     $$ I_1 = \{s \in U_\ast: |g(s)-h_\ast(s)| < \delta\}. $$
     We first estimate the length $|I_1|$.
     Observe that
     $$  h_\ast(0) =  h_\ast(y) = 0 \mbox{ and } |g(0)|=|x-y| < \delta. $$
     Moreover, thanks to Lemma \ref{smalla} and the assumption of this subcase, one has
     \begin{equation}\label{h'}
       |h_\ast'(s)|= |\nabla h(se_1^\ast) \cdot e_1^\ast| \le \frac{a}{10} \mbox{ and } g'(s)= \pm\tan \angle (L,T_y\calC) \ge a, \quad  \forall s\in U_\ast.
     \end{equation}
     (In the following, we assume $g'(s)= \tan \angle (L,T_y\calC)$ and the case $g'(s)= -\tan \angle (L,T_y\calC)$ is similar.)
     Thus
     $$  (g-h_\ast)(s) = (g-h_\ast)(0) + \int_0^s (g'-h_\ast')(\tau) \, d\tau > -\delta + \int_0^s \frac{a}{2} \, d\tau > -\delta + \frac{a}{2}s, \quad s>0. $$
     Thus if $s> \frac{2}{a}\delta$, then $g(s)>h_\ast(s)+\delta$, which implies $(se_1^\ast,g(s)) \notin \calC^\delta$. Similarly one can show that if $s< -\frac{2}{a}\delta$, then $g(s)< h_\ast(s)-\delta$.
     Thus the interval $I_1$ has length $\lesssim_a \delta$.

     We then estimate $|I_2|$.
     Note that for each point $q \in I_2$, we have
     $$  q \in (h_\ast(s_q)-\delta, h_\ast(s_q)+\delta) \mbox{ for some $s_q\in I_1$}.$$
     Thus
     $$  |I_2| \le |h_\ast(I_1)|+2\delta.   $$
      Thanks to \eqref{h'}, we deduce that
      $$  |I_2| \le \max_{s\in I_1}\{|h'_\ast(s)|\}|I_1|+2\delta < \frac{a}{10}|I_1|+2\delta \lesssim_a \delta.   $$
      Finally, since in this subcase $\sin \angle(L,T_y\calC) \sim \angle(L,T_y\calC) \gtrsim a$, thanks to the second inequality in \eqref{I12}, we conclude that
      $$  |I| \sim_a |I_2| \lesssim_a \delta   $$
      as desired.

  %   \medskip
%    \emph{Subcase 1-3. $\angle (L,T_y\calC) \in [ 0, \frac{a}{10})$.}
%     Let the notation be as in Subcase 1-2. Then we have
%     $$  \angle (L,e_1^\ast) = \angle (L,T_y\calC)< \frac{a}{10}.   $$
%     By the standing assumption that $\angle (L, w)  > a$ for all $w \in \calC$, we know
%     $$ \angle (w, e_1^\ast) > \frac{a}{2}.  $$
%     Recalling that the eigenspace of $\nabla^2 h(w)$ with respect to eigenvalue $0$ is span$\{w\}$ and all other eigenvalues are $> \frac{\kappa_1}{2}$, we know
%     $$  h_\ast''(s) \ge \sin \angle (se_1^\ast, e_1^\ast) \frac{\kappa_1}{2} \gtrsim_a \kappa_1, \quad s \in U_\ast. $$
%     Since $g$ is a linear function, we have $g''\equiv 0$ and thus $|g''(s)-h_\ast''(s)| \gtrsim_a \kappa_1$ for all $s \in U_\ast$.
%
%     We claim that
%     $$\|g-h_\ast\|_{C^2(U_\ast)} \sim \inf_{s \in U_\ast}\{|g''(s)-h_\ast''(s)|\}.$$


\medskip
    \emph{Case 2. $I \cap B(o,10 \delta) \ne \emptyset$.}
In this case, since the interval $I \subset L$, we let
$$ I_1':= I \cap B(o,10 \delta) \mbox{ and } I_2':= I \setminus B(o,10 \delta).  $$
Thus
$$  |I| = |I_1'| + |I_2'|. $$
We immediately have $|I_1'| \le \diam B(o,10 \delta) \le 20 \delta.$
In a similar way as in Case 1, one has $|I_2'| \lesssim_a \delta$, which then implies that
$|I| \lesssim_a \delta$ as desired.
The proof is complete.
\end{proof}



\section{Proof of Theorem \ref{cormain}} \label{s5}
In this section, we will first prove Lemma \ref{volume2} and Theorem \ref{cormain}.



Recall the cinematic maps $f_z$ from \eqref{fw},
$$f_z(x):=(e_1(x)\cdot z, \cdots, e_{n-2}(x)\cdot z,  \nu(x)\cdot z) \quad \forall x\in [0,1]^{n-2}$$
where $\{e_i\}_{1 \le i \le n-2}$ is a basis of $T_{\Sigma(x)}\Sigma$ and $\nu(x)\subset T_{\Sigma(x)}S^{n-1}$ is the normal of $\Sigma\subset S^{n-1}$
and the cone with apex $z$ and cross-section $\Sigma$
$$ \calC_z = \bigcup_{x \in \Sigma,r \in [-1,1]} z + rx = \bigcup_{x \in \Sigma,r \in [-1,1]} L_{z,x}. $$
Note that if $y \in L_{w,x}$, then $L_{y,x}$ and $L_{w,x}$ are contained in the same line in $\R^n$.




\begin{lemma} \label{pairs}
    Let $Z \subset B(o,1/2) \subset \R^n$ and $z \in Z$. For any $z' \in Z \setminus\{z\}$,
    if $f_z^\delta \cap f_{z'}^\delta \ne \emptyset$, then $z' \in \calC_w^{2\delta}$.
\end{lemma}


\begin{proof}
   Since $f_z^\delta \cap f_{z'}^\delta \ne \emptyset$, we know that there exists $x \in [0,1]^{n-2}$ such that
   $$ |f_z(x)-f_{z'}(x)| < 2\delta. $$
   This implies that the orthogonal projection $\pi_{P_x}(z)$ and $\pi_{P_{x}}(z')$ of $z$ and of $z'$ to the hyperplane $P_{x}= T_{\Sigma(x)}S^{n-1}$
   satisfies
   \begin{equation}\label{zz'}
     |\pi_{P_{x}}(z)-\pi_{P_{x}}(z')| < 2\delta.
   \end{equation}
   Recalling that $\Sigma(x)$ is the unit normal of $P_{x}$ and $z' \in Z \subset B(o,1/2)$, we have
   $$  z' = \pi_{P_{x}}(z') + r_{z'}\Sigma(x) \mbox{ for some $r_{z'} \in [-1/2,1/2]$}.  $$
   Combining \eqref{zz'}, we deduce that
   \begin{equation}\label{z'z}
     |z' - (\pi_{P_{x}}(z) + r_{z'}\Sigma(x))|  < 2\delta.
   \end{equation}
   Moreover,
   by $z \subset Z \subset B(o,1/2)$, we have $|\pi_{P_{x}}(z)-z| < \frac12$ and hence
   $$  \pi_{P_{x}}(z) + r_{z'}\Sigma(x) \in  \bigcup_{r \in [-1,1]} z + r\Sigma(x) = L_{z,\Sigma(x)} \subset \calC_z.  $$
   Combining \eqref{z'z}, we conclude
   $$   z' \in \calC_z^{2 \delta}$$
   as desired.
\end{proof}


\begin{lemma} \label{Lambda}
    Under the assumption on $\Sigma$ in Theorem \ref{cormain}, for any $y \in S^{n-1}$, define
    $$  \Lambda_y:=  \{x \in \Sigma : \mbox{\rm span}\{y\} \subset T_x\calC_o\}.  $$
    Then
    $$  \dim \Lambda_y \le n-3.  $$
\end{lemma}

\begin{proof}
    Thanks to the proof of Lemma \ref{dual}, we know $T_x\calC_o = T_{\nu(x)} S^{n-1}$ for all $x \in \Sigma$ where we recall $\nu$ is the unit normal of $\Sigma$ at $x$. Thus we deduce that $x \in \Lambda_y$ is equivalent to $\nu(x) \in y^\perp$. Recalling $\nu(\Lambda_y) \subset \nu(\Sigma) = \Sigma^\ast$,
it suffices to check
\begin{equation}\label{n-3}
   \dim (\Sigma^\ast \cap   y^\perp)  \le n-3.
\end{equation}


    By Lemma \ref{dual}, $\Sigma^\ast$ has nonvanishing geodesic curvature ($n=3$) and sectional curvature $> 1$ ($n \ge 4$).
    Since $y^\perp \cap S^{n-1}$ is an $(n-2)$-equatorial sphere which has sectional curvature $\equiv 1$ and $\dim \Sigma^\ast= n-2$, $\Sigma^\ast \cap y^\perp$ can not contain any open subset of $\Sigma^\ast$, which implies that $\Sigma^\ast \cap y^\perp$ has no interior as a subset of $\Sigma^\ast$. As a result,
    $\Sigma^\ast \cap y^\perp$ can not be $(n-2)$-dimensional. Noting that $\Sigma^\ast \cap   y^\perp$ is a submanifold of $S^{n-1}$, $\Sigma^\ast \cap   y^\perp$ has integer dimension. Hence \eqref{n-3} holds.
\end{proof}


We are ready to show Lemma \ref{volume2} and Theorem \ref{cormain}.

\begin{proof}[Proof of Lemma \ref{volume2} and Theorem \ref{cormain}]
    Recall the point $y \in S^{n-1}$ in Theorem \ref{cormain}.
    It suffices to show that Theorem \ref{cormain} holds for $\Sigma \setminus (\cup_{x \in \Lambda_y}B(x,a))$ for each $a>0$ small enough where $\Lambda_y$ is the set in Lemma \ref{Lambda}.
    Thus we fix $a>0$ small enough and partition $\Sigma \setminus (\cup_{x \in \Lambda_y}B(x,a))$ into small pieces such that all the results in Section \ref{s4} hold.
    We prove Theorem \ref{cormain} for each piece. Fix a piece and still denote it by $\Sigma$. In particular, for each $z \in \R^n$ and $p \in \calC_z\setminus\{z\}$, noting that $\frac{p-z}{|p-z|} \in \Sigma$, we have
    \begin{equation}\label{ang2}
      \angle (y, T_p\calC_z)=\angle (y, T_{p-z}\calC_{o})=\angle (y, T_{\frac{p-z}{|p-z|}}\calC_o) \ge \frac{a}2.
    \end{equation}


    We first prove Lemma \ref{volume2}. By Lemma \ref{cineprop}, let $D>1$ be the bilipschitz constant of the bilipschitz map from $Z$ to $F=\{f_{z} \in C^2( [0,1]^{n-2},[0,1]^{n-1})\}_{z \in Z}$. For $\epsilon>0$, we choose $\delta_0=\delta_0(\epsilon)>0$ such that
    $$  \delta_0^\epsilon < \min\{D^{-(n-2)}, \log(1/\delta_0)^{-1}\}.  $$
    We will show Lemma \ref{volume2} holds for all $\delta<\delta_0$.
    Recalling $\dim Z = t$,
     we can find a $D^{-1}\delta$-separated set $Y\subset Z$ with $|Y| \sim \delta^{-t}$. Recall from \eqref{vertical} that $f_z^\delta$ is the vertical $\delta$-neighbourhood of $\mbox{graph}(f_z)$, we deduce that
     $$  \bigcup_{z \in Z}\mbox{graph}(f_z) \subset \bigcup_{z \in Y}f_z^\delta.$$
     Recalling from \eqref{heu31} that
     $$\mathcal{Z} = \bigcup_{z \in Z} \bigcup_{x \in X_z}(z, f_z(z)) \subset [0,1]^{2n-3}$$
     where
     $X_z \subset [0,1]^{n-2}$ with $\calL^{n-2}(X_z)>C_Z>0$ for each $z \in Z$. By \cite[Lemma 2.7]{orponen2023hausdorff}, we can find a subset $\delta$-separated $(\delta,n-2,\delta^{-\epsilon})$-set $\widetilde{X}_z \subset X_z$ with $\delta^{\epsilon}\delta^{-(n-2)} \le |\widetilde{X}_z| \le \delta^{-(n-2)}$.
    To show
    $ |\mathcal{Z}|_{\delta} \gtrsim \delta^{3\epsilon - (n-2 +t)}$ in Lemma \ref{volume2},
    it suffices to show
    \begin{equation}\label{caly}
      |\mathcal{Y}|_{\delta}\gtrsim \delta^{3\epsilon - (n-2 +t)}
    \end{equation}
    where
     $$  \mathcal{Y} := \bigcup_{z \in Y} \bigcup_{x \in \widetilde{X}_z }(x, f_z(x)). $$
The following proof is devoted to showing \eqref{caly}.

Recalling from \eqref{le1} that $y \in S^{n-1}$ is the point such that
    $$ \dim \pi_{T_yS^{n-1}}(Z) \le n-2.$$
By \cite[Lemma 2.7]{orponen2023hausdorff}, for any $s < \dim \pi_{T_yS^{n-1}}(Z)$, we can find a $\delta$-separated $(\delta,s,\delta^{-\epsilon})$-set $W \subset \pi_{T_yS^{n-1}}(Z)$ with $|W| \le \delta^{-s}$. Moreover, $\pi_{T_yS^{n-1}}(Z) \subset W^{n\delta}$ where $W^{n\delta}= \cup_{w \in W} B(w, n \delta)$ is the $n\delta$-neighbourhood of $W$ and thus the preimage set satisfies
\begin{equation}\label{ndz}
  \bigcup_{w \in W}\pi_{T_yS^{n-1}}^{-1}(B(w,n\delta)) = \pi_{T_yS^{n-1}}^{-1}(W^{n\delta})  \supset Z \supset Y.
\end{equation}
Recalling that $Y \subset B(o,1/2)$, we observe that $\pi_{T_yS^{n-1}}^{-1}(B(w,n\delta))\cap Y$ is contained in the $n\delta$-neighbourhood $L_{w,y}^{n\delta}$ of the line segment $L_{w,y}= \cup_{r \in [-1,1]} \{w + ry\}$.
Combining \eqref{ndz}, we deduce that
\begin{equation}\label{incly}
  Y \subset   \bigcup_{w \in W}L_{w,y}^{n\delta}.
\end{equation}
%Since $W$ is $\delta$-separated, $\cup_{w \in W} B(w, n \delta)$ has $C(n)$-overlapping for some $C(n)>0$.
%For each $w \in W$, define
%$$W_w:=\pi_{T_yS^{n-1}}^{-1}(B(w,n\delta)) \cap Y.$$
%We know $W_w$ is contained in the $n\delta$-neighbourhood $L_{w,y}^{n\delta}$ of the line segment $L_{w,y}$. In particular,
%  $$W_w \subset \calC_w^{n\delta} \subset L_{w,y}^{n\delta}.$$
Noting that $\{L_{w,y}\}_{w \in W}$ are parallel line segments perpendicular to $T_yS^{n-1}$, we have
$$ \pi_{T_yS^{n-1}}(L_{w,y}) = w \mbox{ and } \pi_{T_yS^{n-1}}(L_{w,y}^{n\delta}) = B^{n-1}(w,n\delta) \subset T_yS^{n-1}. $$
Thus for any ball $B^n(x,r) \subset \R^n$ with $r \in [\delta,1]$, we have
\begin{align*}
  |\{w \in W: L_{w,y}^{n\delta} \cap B^n(x,r) \ne \emptyset\}| & = |\{w \in W: \pi_{T_yS^{n-1}}[L_{w,y}^{n\delta} \cap B^n(x,r)] \ne \emptyset\}| \\
   & =|\{w \in W: B^{n-1}(w,n\delta) \cap B^{n-1}(\pi_{T_yS^{n-1}}(x),r) \ne \emptyset\}| \\
   & \lesssim |W^{n\delta} \cap B^{n-1}(\pi_{T_yS^{n-1}}(x),r)|_{\delta} \\
   & \lesssim |W \cap B^{n-1}(\pi_{T_yS^{n-1}}(x),r)|_{\delta} \le \delta^{-\epsilon}r^s|W|
\end{align*}
where in the last inequality we recall that $W$ is a $(\delta,s,\delta^{-\epsilon})$-set.

For each $z \in Y$, define
$$  Y_z : = \calC_z^{n\delta} \cap (Y\setminus \{z\}).  $$
For each $r \in [\delta,1]$, thanks to \eqref{incly}, recalling \eqref{ang2} and applying Theorem \ref{lineconen}, we obtain that
\begin{equation}\label{ds}
    |Y_z \cap B(x,r)|  \le  |\{w \in W: L_{w,y}^{n\delta} \cap B(x,r) \ne \emptyset\}| \cdot |L_{w,y}^{n\delta} \cap \calC_z^{n\delta}|_\delta \lesssim_a \delta^{-\epsilon}r^s|W|.
\end{equation}
   For each dyadic $r \in [\delta,1]$ and $z \in Y$, define
   $$  Y_{z,r}:= \{ z' \in Y_z : |z'-z|  \in (r,2r]  \} . $$
   We deduce from \eqref{ds} that
   \begin{equation}\label{yzr}
     |Y_{z,r}| \le |Y_z \cap B(z,2r)| \lesssim_a \delta^{-\epsilon}r^s|W| \le \delta^{-\epsilon}r^s\delta^{-s} .
   \end{equation}


   Moreover, recall from Lemma \ref{cineprop} that
  $F=\{f_{z} \in C^2( [0,1]^{n-2},[0,1]^{n-1})\}_{z \in Z}$ is a cinematic family and $F$ is biLipschitz homeomorphic to $Z$. Applying Lemma \ref{enbhd}, we have
  \begin{equation}\label{bilip}
    \mathcal{L}^{2n-3}( f^\delta_z(X_z^\delta)\cap  f^\delta_{z'}(X_{z'}^\delta)) \lesssim \frac{\delta^{2n-3}}{\|f_z-f_{z'}\|_{C^2([0,1]^{n-2})}^{n-2}} \le D^{n-2} \frac{\delta^{2n-3}}{|z-z'|^{n-2}} \le \frac{\delta^{-\epsilon}\delta^{2n-3}}{|z-z'|^{n-2}}.
  \end{equation}
  Combining \eqref{yzr} and \eqref{bilip}, we arrive at
    \begin{align*}
\sum_{z,z'\in Y}\mathcal{L}^{2n-3}( f^\delta_z(X_z^\delta)\cap  f^\delta_{z'}(X_{z'}^\delta)) &= \sum_{z\in Y}\sum_{r=\delta}^1\sum_{z'\in Y_{z,r}} \mathcal{L}^{2n-3}( f^\delta_z(X_z^\delta)\cap  f^\delta_{z'}(X_{z'}^\delta))  \\
& \lesssim \sum_{z\in Y}\sum_{r=\delta}^1 \sum_{z'\in Y_{z,r}} \frac{\delta^{2n-3}}{\|f_z-f_{z'}\|_{C^2([0,1]^{n-2})}^{n-2}}
 \le \sum_{z\in Y}\sum_{r=\delta}^1 \sum_{z'\in Y_{z,r}} \frac{\delta^{-\epsilon}\delta^{2n-3}}{|z-z'|^{n-2}} \\
& \le \sum_{z\in Y}\sum_{r=\delta}^1 |Y_{z,r}|\frac{\delta^{-\epsilon}\delta^{2n-3}}{r^{n-2}}
 \le \sum_{z\in Y}\sum_{r=\delta}^1 \delta^{-2\epsilon} r^{s}\delta^{-s}\frac{\delta^{2n-3}}{r^{n-2}} \\
 & = |Y|\log(1/\delta)\delta^{-2\epsilon} \frac{\delta^{2n-3-s}}{r^{n-2-s}} \le \delta^{-3\epsilon}
  \delta^{n-1}|Y|
  \end{align*}
where in the last line we recall $s \le n-2$ from the assumption \eqref{le1}.
Similar to inequality \eqref{cs} in Section \ref{s3}, we conclude that
 \begin{align*}
  \mathcal{L}^{2n-3}(\mathcal{Y}) &=  \mathcal{L}^{2n-3}(\bigcup_{z \in Y}  f^\delta_z(X_z^\delta))  \ge \frac{[\sum_{z \in Y}\mathcal{L}^{2n-3}( f^\delta_z(X_z^\delta)]^2}{\sum_{z,z'\in Y}\mathcal{L}^{2n-3}( f^\delta_z(X_z^\delta)\cap  f^\delta_{z'}(X_{z'}^\delta))}  \\
  & \gtrsim \frac{|Y|^2 \delta^{2n-2}}{\delta^{-3\epsilon} \delta^{n-1}|Y|} = \delta^{3\epsilon} \delta^{n-1}|Y|.
  \end{align*}
  Recalling $|Y| \gtrsim \delta^{-t}$, we obtain \eqref{caly}
as desired. The proof of Lemma \ref{volume2} is complete.

Finally, Theorem \ref{cormain} can be shown in the same way as the proof of Theorem \ref{planemain} in Section \ref{s3} with the help of Lemma \ref{volume2}. We omit the details.
\end{proof}





\bibliographystyle{plain}
\bibliography{references.bib}



\end{document}
